\documentclass[UTF8,fontset=none,11pt]{ctexart}

\usepackage[a4paper,margin=2.2cm]{geometry}
\usepackage{amsmath,amssymb,amsthm,mathtools,bm}
\usepackage{booktabs}
\usepackage{enumitem}
\usepackage{xcolor}
\usepackage{hyperref}
\usepackage{microtype}
\usepackage{array}

\setCJKmainfont{Noto Serif CJK SC}
\setCJKsansfont{Noto Sans CJK SC}
\setCJKmonofont{Noto Sans CJK SC}
\setmainfont{DejaVu Serif}
\setsansfont{DejaVu Sans}
\setmonofont{DejaVu Sans Mono}

\hypersetup{
  colorlinks=true,
  linkcolor=blue!55!black,
  citecolor=blue!55!black,
  urlcolor=blue!55!black
}

\allowdisplaybreaks
\setlist[itemize]{leftmargin=1.8em,itemsep=0.25em,topsep=0.25em}
\setlist[enumerate]{leftmargin=2.0em,itemsep=0.25em,topsep=0.25em}
\renewcommand{\arraystretch}{1.18}

\newtheoremstyle{cnstyle}
  {0.6em}{0.6em}{\normalfont}{}{\bfseries}{.}{0.5em}{}
\theoremstyle{cnstyle}
\newtheorem{theorem}{定理}[section]
\newtheorem{lemma}[theorem]{引理}
\newtheorem{proposition}[theorem]{命题}
\newtheorem{assumption}[theorem]{假设}
\newtheorem{definition}[theorem]{定义}
\newtheorem{remark}[theorem]{说明}

\newcommand{\R}{\mathbb{R}}
\newcommand{\E}{\mathbb{E}}
\newcommand{\Cov}{\operatorname{Cov}}
\newcommand{\tr}{\operatorname{tr}}
\newcommand{\rank}{\operatorname{rank}}
\newcommand{\diag}{\operatorname{diag}}
\newcommand{\vect}{\operatorname{vec}}
\newcommand{\argmin}{\operatorname*{arg\,min}}
\newcommand{\argmax}{\operatorname*{arg\,max}}
\newcommand{\calR}{\mathcal{R}}
\newcommand{\calL}{\mathcal{L}}
\newcommand{\calD}{\mathcal{D}}
\newcommand{\calI}{\mathcal{I}}
\newcommand{\ip}[2]{\left\langle #1,#2\right\rangle}
\newcommand{\norm}[1]{\left\lVert #1\right\rVert}
\newcommand{\fnorm}[1]{\left\lVert #1\right\rVert_F}
\newcommand{\twonorm}[1]{\left\lVert #1\right\rVert_2}
\newcommand{\psinv}{^{+}}
\newcommand{\st}{\;\mathrm{s.t.}\;}
\newcommand{\dd}{\,\mathrm{d}}

\title{\bfseries 论文公式与证明导读\\
{\large Theory of Continual Learning Against Data Poisoning Attacks}}
\author{阅读笔记与证明讲解}
\date{\today}

\begin{document}
\maketitle
\tableofcontents
\newpage

\section{先给结论：这篇论文到底在证明什么}

这篇论文研究的是：在连续学习（continual learning, CL）中，模型按任务序列
\(1,2,\dots,T\) 依次训练。如果攻击者在部分任务的数据里加入投毒扰动，
正则化式连续学习方法什么时候一定防不住，什么时候可以设计出有证明保证的防御。

论文的三条主线如下。

\begin{enumerate}
  \item \textbf{不可防御边界。}
  如果攻击很频繁，并且攻击能改变数据均值（shifted attack），或者虽然不改变均值但每次攻击方差预算随 \(T\) 线性增长，那么任何形如
  \[
    w_t\in\argmin_w
    \left\{
      \frac{1}{n_t}\sum_{z\in D_t}\ell(w,z)
      +\frac12\norm{w-w_{t-1}}_{H_t}^2
    \right\}
  \]
  的正则化式 CL 方法都无法保证最终误差收敛到 0。

  \item \textbf{少量但可能无限强的攻击可以用检测加回滚防御。}
  对于只攻击有限多个任务但单次攻击强度可以很大的情形，论文设计了 task-to-task verification（T2T）检测分数 \(d_t\)。它通过两个相邻更新
  \(w_t-w_{t-1}\) 与 \(w_{t-1}-w_{t-2}\) 的投影组合，消去历史状态
  \(w_{t-2}\) 的影响。最终证明：只要攻击任务数 \(K\) 不随 \(T\) 线性增长，最终误差可达到
  \[
    \calR(w_T)=O\!\left(\frac{K^2}{\epsilon T}\right)
  \]
  量级。

  \item \textbf{频繁但有界、零均值攻击需要鲁棒特征正则化。}
  如果几乎每个任务都可能被投毒，但攻击是 non-shifted 且有界，单纯检测会丢掉太多任务。论文把防御写成在线 minimax 问题，并在可对角化假设下给出每个特征方向的最优正则化强度。结论是：鲁棒正则化不让攻击集中打最脆弱的单一方向，而是把风险分摊到一组 protected directions。
\end{enumerate}

\section{背景知识补充}

\subsection{连续学习与灾难性遗忘}

连续学习的特点是数据不是一次性给出，而是按任务序列到达。第 \(t\) 个任务有数据集
\[
  D_t=\{(x_t^i,y_t^i)\}_{i=1}^{n_t},\qquad
  x_t^i\in\mathcal X,\; y_t^i\in\mathcal Y .
\]
训练第 \(t\) 个任务时，通常不能再完整访问旧任务数据。因此模型既要学习新任务，又不能忘掉旧任务，这就是稳定性-可塑性权衡。

论文用全局超额风险衡量最终模型：
\begin{equation}
  \calR_T(w)
  =
  \sum_{t=1}^T \E_{z\sim\calD_t}[\ell(w,z)]
  -
  \min_{w^\star}\sum_{t=1}^T\E_{z\sim\calD_t}[\ell(w^\star,z)] .
  \label{eq:global-risk}
\end{equation}
它问的是：当前模型 \(w\) 和最优全局模型相比，多损失了多少。

\subsection{第一个风险公式逐符号解释}

公式 \eqref{eq:global-risk} 是论文最早出现、也最容易让人迷路的公式。把它拆开看：
\[
  \calR_T(w)
  =
  \underbrace{\sum_{t=1}^T \E_{z\sim\calD_t}[\ell(w,z)]}_{\text{当前模型 }w\text{ 在所有任务上的总期望损失}}
  -
  \underbrace{\min_{u}\sum_{t=1}^T\E_{z\sim\calD_t}[\ell(u,z)]}_{\text{所有模型里能达到的最小总期望损失}} .
\]
这里每个符号的含义是：
\begin{itemize}
  \item \(T\)：总任务数。
  \item \(t\)：任务下标，\(t=1,\dots,T\)。
  \item \(\calD_t\)：第 \(t\) 个任务背后的真实数据分布。它不是数据集本身，而是产生数据的概率规律。
  \item \(z\)：一个样本的简写。全文默认 \(z=(x,y)\)，其中 \(x\) 是输入特征，\(y\) 是标签或响应。
  \item \(z\sim\calD_t\)：从第 \(t\) 个任务的数据分布中随机抽一个样本。
  \item \(\ell(w,z)\)：模型参数为 \(w\) 时，在样本 \(z=(x,y)\) 上的损失。例如平方损失是
  \[
    \ell(w,(x,y))=\frac12\norm{x^\top w-y}_2^2 .
  \]
  \item \(\E_{z\sim\calD_t}[\ell(w,z)]\)：模型 \(w\) 在第 \(t\) 个任务上的平均损失。注意这是对真实分布求平均，不只是对训练集求平均。
  \item \(\sum_{t=1}^T\)：把所有任务的损失加起来。
  \item \(\min_u\)：在所有可能参数 \(u\) 里找总损失最小的那个值。这里用 \(u\) 是为了避免和真正选出的最优模型混淆。
\end{itemize}

所谓\textbf{最优全局模型}就是
\[
  w_{\mathrm{glob}}^\star
  \in
  \argmin_u
  \sum_{t=1}^T\E_{z\sim\calD_t}[\ell(u,z)] .
\]
它是“如果你事先知道所有任务的真实分布，并且可以一次性训练，那么最好的参数”。CL 算法实际只能顺序看到任务，所以 \(w_T\) 通常达不到 \(w_{\mathrm{glob}}^\star\)。超额风险 \(\calR_T(w_T)\) 就是在衡量差距。

如果最优解不唯一，理论里常见做法是任选一个最优解，或者在所有最优解里选范数最小的那个。论文后面的线性回归设置通常直接假设存在共享真实参数 \(w^\star\)，这时 \(w^\star\) 就扮演全局目标模型。

\subsection{经验损失、\(Z_t\)、\(\calL(w_t,Z_\tau)\)}

真实期望 \(\E_{z\sim\calD_t}\) 通常算不出来，因为我们不知道 \(\calD_t\)。训练时能用的是数据集
\[
  D_t=\{z_t^1,\dots,z_t^{n_t}\},
  \qquad
  z_t^i=(x_t^i,y_t^i).
\]
论文有时用 \(Z_t\) 表示“把任务 \(t\) 的全部样本拼起来形成的数据矩阵/数据集合”。你可以先把 \(Z_t\) 理解成 \(D_t\) 的矩阵版。

任务 \(t\) 的\textbf{经验损失}定义为
\[
  \calL(w,Z_t)
  =
  \frac1{n_t}\sum_{i=1}^{n_t}\ell(w,z_t^i).
\]
所以
\[
  \calL(w_t,Z_\tau)
  =
  \frac1{n_\tau}\sum_{i=1}^{n_\tau}\ell(w_t,z_\tau^i)
\]
的意思是：把“第 \(t\) 步得到的模型 \(w_t\)”拿去评估“第 \(\tau\) 个任务的数据 \(Z_\tau\)”上的平均损失。

论文中出现的经验损失和
\[
  \sum_{\tau=1}^t\calL(w_t,Z_\tau)
\]
就是
\[
  \calL(w_t,Z_1)+\calL(w_t,Z_2)+\cdots+\calL(w_t,Z_t),
\]
意思是：当前模型 \(w_t\) 在已经见过的所有任务训练集上的总经验损失。它是下面这个真实全局风险的可计算替代物：
\[
  \sum_{\tau=1}^t\E_{z\sim\calD_\tau}[\ell(w_t,z)].
\]

\subsection{正则化式 CL 更新}

论文研究的核心算法族是
\begin{equation}
  w_t\in\argmin_{w'}
  \left\{
  \frac{1}{n_t}\sum_{z\in D_t}\ell(w',z)
  +\frac12\norm{w'-w_{t-1}}_{H_t}^2
  \right\},
  \label{eq:regularized-cl}
\end{equation}
其中
\[
  \norm{u}_{H_t}^2=u^\top H_tu .
\]
第一项要求模型适配当前任务；第二项惩罚离开上一轮模型 \(w_{t-1}\) 太远。
矩阵 \(H_t\succeq 0\) 决定哪些参数方向更该被保护。

为什么“适配当前任务”对应第一项？如果完全不考虑旧任务，第 \(t\) 步最自然的训练目标就是最小化当前任务训练集上的平均损失：
\[
  \min_{w'}\frac1{n_t}\sum_{z\in D_t}\ell(w',z).
\]
这就是普通监督学习的经验风险最小化。CL 的问题是不能只管当前任务，否则会忘掉过去。因此论文在这个目标后面加上
\[
  \frac12\norm{w'-w_{t-1}}_{H_t}^2,
\]
要求新参数 \(w'\) 不要离上一轮参数 \(w_{t-1}\) 太远。

为什么要引入 \(H_t\)？如果只用一个标量 \(\lambda\norm{w'-w_{t-1}}_2^2\)，所有参数方向被同等保护。但真实模型里有些方向对旧任务很重要，有些方向不重要。矩阵 \(H_t\) 可以给不同方向不同惩罚强度。若 \(H_t\) 在某个方向上特征值大，模型就很难沿那个方向移动；若特征值小，模型可以更自由地更新。

直觉上：
\begin{itemize}
  \item \(H_t=0\)：只管新任务，容易忘旧任务。
  \item \(H_t\) 很大：几乎冻结模型，旧知识保住了，但学不动新任务。
  \item \(H_t\) 是矩阵而不是标量：不同特征方向可以有不同保护强度。
\end{itemize}

\subsection{范数符号：两个竖线和下标是什么意思}

论文里大量出现 \(\norm{\cdot}\)、\(\norm{\cdot}_2\)、\(\fnorm{\cdot}\)、\(\norm{\cdot}_{H}\)。它们都表示“大小”，但对象不同。

\begin{itemize}
  \item 向量二范数：
  \[
    \norm{v}_2=\sqrt{v_1^2+\cdots+v_d^2}.
  \]
  下标 \(2\) 表示平方和再开方。没有下标时，很多论文默认也是二范数。

  \item 矩阵 Frobenius 范数：
  \[
    \norm{A}_F=\sqrt{\sum_{i,j}A_{ij}^2}.
  \]
  它就是把矩阵所有元素平方求和再开方。

  \item 矩阵谱范数：
  \[
    \norm{A}_2=\max_{\norm{x}_2=1}\norm{Ax}_2.
  \]
  它表示矩阵最多能把一个单位向量放大多少倍，也等于最大奇异值。

  \item 加权范数：
  \[
    \norm{u}_{H}^2=u^\top Hu.
  \]
  若 \(H=I\)，它就是普通二范数平方；若 \(H\neq I\)，不同方向被赋予不同权重。
\end{itemize}

\subsection{数据投毒的三个维度}

攻击者在任务 \(t\) 上加入扰动
\[
  \eta_t=(\eta_{x,t},\eta_{y,t}),\qquad
  \E[\fnorm{\eta_t}^2]\le M .
\]
论文把扰动分解成
\[
  \eta_t=\hat\eta_t+\tilde\eta_t ,
\]
其中 \(\hat\eta_t\) 是条件均值偏移，\(\tilde\eta_t\) 是零均值随机扰动：
\[
  \E[\tilde\eta_t\mid\calI_t]=0 .
\]
为什么扰动既写在 \(x\) 又写在 \(y\)？因为训练数据样本是 \((x,y)\)。现实攻击可以改输入，例如图像里加触发器、传感器特征被污染；也可以改标签，例如把类别标错或在线性回归中改变响应值。论文先给出最一般的投毒形式
\[
  (x_t^i,y_t^i)\mapsto(x_t^i+\eta_{x,t}^i,\;y_t^i+\eta_{y,t}^i),
\]
后面为了可证明性，在线性回归理论分析中主要研究 label poisoning，也就是只改 \(y\)，不改 \(x\)。

零均值随机扰动
\[
  \E[\tilde\eta_t\mid\calI_t]=0
\]
的意思是：在攻击者已经知道历史、当前模型、当前数据等信息 \(\calI_t\) 之后，\(\tilde\eta_t\) 的条件平均仍然是 0。它不持续把数据往某个固定方向推偏，但可以增加方差。举例：\(\tilde\eta_t\) 以一半概率取 \(+a\)，一半概率取 \(-a\)，平均就是 0。

三个攻击维度是：

\begin{center}
\begin{tabular}{lll}
\toprule
维度 & 类型 & 含义 \\
\midrule
频率 & frequent / infrequent & 攻击任务数 \(K=\Theta(T)\) 或 \(K=O(1)\) \\
强度 & unbounded / bounded & 单次预算 \(M=\Omega(T)\) 或 \(M=o(T)\) \\
模式 & shifted / non-shifted & 是否有均值偏移 \(\hat\eta_t\neq 0\) \\
\bottomrule
\end{tabular}
\end{center}

\subsection{在线 minimax 问题是什么}

minimax 的字面意思是“我先选策略，使得对手在最坏情况下造成的损失最小”。形式上常写成
\[
  \min_{\text{防御策略}}\max_{\text{攻击策略}}\;\text{损失}.
\]
在这篇论文里，每个时刻 \(t\) 都有一次对抗：
\[
  \min_{H_t}\max_{\eta_t}\calR_t(w_t),
\]
防御者选正则化矩阵 \(H_t\)，攻击者选投毒扰动 \(\eta_t\)，双方围绕最终风险对抗。

为什么叫\textbf{在线}？因为任务按时间到达。第 \(t\) 轮只能用当前和历史信息，不能提前知道未来任务。在线 minimax 通常有几类解决方法：
\begin{itemize}
  \item 直接求 saddle point，即找一组策略让双方都无法单方面改进。
  \item 把内层最大化先解掉，变成一个单独的最小化问题。
  \item 用凸优化和 KKT 条件求闭式解。
  \item 在矩阵问题太复杂时，通过对角化把矩阵问题拆成多个一维标量问题。
\end{itemize}

这篇论文为什么需要对角化？因为 \(H_t\)、\(X_t^\top X_t\)、协方差矩阵都是高维矩阵，直接求解 minimax 很难。如果假设这些 Hessian 可以同时对角化，那么每个特征方向 \(u_j\) 互不耦合，整个问题就变成“给每个方向选一个标量 \(\lambda_j^t\)”的问题。这样才能用 KKT 条件推出 Proposition 5.2 的闭式水位解。

\subsection{Taylor 展开、一阶最优性和曲率}

一维 Taylor 展开是
\[
  f(a+h)
  =
  f(a)+f'(a)h+\frac12f''(a)h^2+O(h^3).
\]
多维情形中，若 \(w\) 是向量，在点 \(w_0\) 附近展开：
\[
  f(w)
  =
  f(w_0)
  +\nabla f(w_0)^\top(w-w_0)
  +\frac12(w-w_0)^\top\nabla^2f(w_0)(w-w_0)
  +O(\norm{w-w_0}^3).
\]
这里：
\begin{itemize}
  \item \(\nabla f(w_0)\) 是梯度，一阶导数向量。
  \item \(\nabla^2 f(w_0)\) 是 Hessian，二阶导数组成的矩阵。
  \item 曲率通常就指二阶导数或 Hessian。曲率大表示函数在这个方向变化很快，参数稍微动一点损失就变很多。
\end{itemize}

一阶最优性条件是：如果可微函数 \(f(w)\) 在内部点 \(w^\star\) 取得局部最小值，那么
\[
  \nabla f(w^\star)=0.
\]
直观上，最小点附近不能再有“下降方向”，所以所有一阶方向导数都为 0。论文里每次写
\[
  \nabla_w\calL(w_t,Z_t+\eta_t)+H_t(w_t-w_{t-1})=0
\]
就是对正则化目标
\[
  \calL(w,Z_t+\eta_t)+\frac12\norm{w-w_{t-1}}_{H_t}^2
\]
求梯度并令其为 0。第二项的梯度是
\[
  \nabla_w\frac12(w-w_{t-1})^\top H_t(w-w_{t-1})
  =
  H_t(w-w_{t-1}),
\]
这里使用了 \(H_t\) 对称这一事实。

\subsection{为什么会出现 Hessian、伪逆和混合二阶导}

很多公式看起来复杂，是因为论文把“数据扰动如何影响参数更新”局部线性化。
若
\[
  Q_t=\nabla_{ww}^2\calL(w_t^\star,Z_t),
  \qquad
  G_t=\nabla_{wZ}^2\calL(w_t^\star,Z_t),
  \qquad
  S_t=Q_t+H_t,
\]
那么：
\begin{itemize}
  \item \(Q_t\) 是当前任务损失对参数的曲率；
  \item \(G_t\) 衡量数据 \(Z_t\) 变化如何改变参数梯度；
  \item \(S_t=Q_t+H_t\) 是“数据曲率 + 正则化曲率”；
  \item \(S_t^+\) 是 Moore-Penrose 伪逆，用来处理 \(S_t\) 不满秩的情况。
\end{itemize}

局部看，投毒先通过 \(G_t\eta_t\) 改变梯度，再通过 \(S_t^+\) 转化为参数偏移。
所以敏感矩阵经常长成
\[
  S_t^+G_t
  =
  (\nabla_{ww}^2\calL_t+H_t)^+
  \nabla_{wZ}^2\calL_t .
\]

这里的“敏感矩阵”指的是：输入一个数据扰动 \(\eta_t\)，它输出由此造成的参数偏移或风险偏移。若这个矩阵在某个方向上放大倍数很大，攻击者就会优先沿这个方向投毒。

为什么是“数据曲率 + 正则化曲率” \(S_t=Q_t+H_t\)？因为总训练目标是
\[
  \calL(w,Z_t)+\frac12\norm{w-w_{t-1}}_{H_t}^2 .
\]
第一项对 \(w\) 的 Hessian 是 \(Q_t\)，第二项对 \(w\) 的 Hessian 是 \(H_t\)。两者相加就是总目标对参数的局部二阶变化：
\[
  \nabla_{ww}^2
  \left[
  \calL(w,Z_t)+\frac12\norm{w-w_{t-1}}_{H_t}^2
  \right]
  =
  Q_t+H_t.
\]

为什么 \(G_t=\nabla_{wZ}^2\calL\) 是“先对 \(w\) 再对 \(Z\)”？因为最优性条件首先是关于参数梯度的：
\[
  \nabla_w\calL(w,Z)=0.
\]
当数据从 \(Z_t\) 变成 \(Z_t+\eta_t\) 时，真正受影响的是这个参数梯度。于是我们看“参数梯度 \(\nabla_w\calL\) 对数据 \(Z\) 的导数”，这就是混合导数
\[
  \nabla_Z(\nabla_w\calL)=\nabla_{wZ}^2\calL.
\]
在足够光滑时，先对 \(w\) 再对 \(Z\) 或先对 \(Z\) 再对 \(w\) 本质上是同一个混合二阶信息；论文的写法强调它是“数据扰动如何改变参数梯度”。

\subsection{Moore--Penrose 伪逆是什么}

普通逆矩阵 \(A^{-1}\) 只有在 \(A\) 是方阵且满秩时才存在。但机器学习里常有过参数化模型，Hessian 可能奇异或不满秩。这时方程
\[
  Ax=b
\]
可能没有唯一解，甚至没有精确解。Moore--Penrose 伪逆 \(A^+\) 给出一个统一处理方式：
\begin{itemize}
  \item 如果 \(A\) 可逆，那么 \(A^+=A^{-1}\)。
  \item 如果 \(Ax=b\) 有很多解，\(A^+b\) 是范数最小的解。
  \item 如果没有精确解，\(A^+b\) 给出最小二乘意义下的解。
\end{itemize}

论文里出现 \(S_t^+\) 是因为线性化后要解
\[
  S_t(w_t-w_t^\star)
  \approx
  \text{某个由历史项和投毒项组成的向量}.
\]
若 \(S_t\) 不可逆，就不能写 \(S_t^{-1}\)，只能写更一般的 \(S_t^+\)。

为什么说“投毒先影响梯度，再通过伪逆变成参数偏移”？因为训练后的参数由一阶最优性条件决定。数据一变，梯度方程的右边先变；为了重新满足“梯度为 0”，参数必须移动。局部线性化后这个关系就是
\[
  S_t\cdot(\text{参数移动})
  +
  G_t\cdot(\text{数据扰动})
  \approx 0,
\]
所以
\[
  \text{参数移动}
  \approx
  -S_t^+G_t\cdot(\text{数据扰动}).
\]

\subsection{Kronecker 积、\(I_{p_2}\) 和最大奇异方向}

\(I_{p_2}\) 是 \(p_2\times p_2\) 的单位矩阵。例如 \(p_2=3\) 时，
\[
  I_3=
  \begin{bmatrix}
  1&0&0\\
  0&1&0\\
  0&0&1
  \end{bmatrix}.
\]

Kronecker 积 \(\otimes\) 是一种把两个矩阵扩展成大矩阵的运算。若
\[
  A=
  \begin{bmatrix}
  a&b\\
  c&d
  \end{bmatrix},
  \qquad
  B=
  \begin{bmatrix}
  1&2\\
  3&4
  \end{bmatrix},
\]
则
\[
  A\otimes B=
  \begin{bmatrix}
  aB&bB\\
  cB&dB
  \end{bmatrix}.
\]
在线性回归里，如果输出 \(y\) 是 \(p_2\) 维向量，参数 \(w\) 是矩阵。为了使用普通向量公式，论文把 \(w\) 拉直成 \(\vect(w)\)，并用
\[
  \mathbf X_t=I_{p_2}\otimes X_t
\]
构造一个大设计矩阵，使得
\[
  \vect(Y_t)=\mathbf X_t\vect(w^\star)+\vect(\varepsilon_t).
\]

任意矩阵 \(A\) 都可以做奇异值分解
\[
  A=U\Sigma V^\top,
\]
其中 \(\Sigma\) 对角线上的非负数叫奇异值：
\[
  \sigma_1\ge\sigma_2\ge\cdots\ge0.
\]
最大奇异值 \(\sigma_1\) 表示
\[
  \max_{\norm{x}_2=1}\norm{Ax}_2.
\]
达到这个最大值的输入方向 \(v_1\) 就是最大右奇异方向。攻击者若有预算
\(\E\norm{\eta}^2\le M\)，而目标近似为 \(\E\norm{A\eta}^2\)，那么把扰动放在 \(v_1\) 方向最划算，因为它被 \(A\) 放大最多。

\subsection{线性回归设置}

为了得到闭式证明，论文后半部分使用多输出线性回归：
\begin{equation}
  Y_t=X_tw^\star+\varepsilon_t,
  \qquad
  \E[\varepsilon_t^i\mid x_t^i]=0,
  \qquad
  \E[\varepsilon_t^i(\varepsilon_t^i)^\top\mid x_t^i]=\sigma^2I .
  \label{eq:linear-model}
\end{equation}
平方损失是
\[
  \ell(w,(x,y))=\frac12\norm{x^\top w-y}_2^2 .
\]
论文在这个设置下把参数矩阵向量化：
\[
  \mathbf w_t=\vect(w_t),\qquad
  \mathbf X_t=I_{p_2}\otimes X_t .
\]
这里的 \(\otimes\) 是 Kronecker 积。它的作用是把多输出线性回归写成普通向量形式：
\[
  \vect(Y_t)=\mathbf X_t\vect(w^\star)+\vect(\varepsilon_t).
\]

\section{进一步答疑：把概念和例子补齐}

\subsection{顺序到来的任务里，怎么谈全局最优}

在 CL 过程中，算法通常\textbf{不能真正求出最终全局最优}，因为它在第 \(t\) 步只看到了 \(1,\dots,t\) 的任务，还不知道未来任务 \(t+1,\dots,T\)。论文里的全局最优更多是一个\textbf{评价基准}：
\[
  w_{\mathrm{glob},T}^\star
  \in
  \argmin_w
  \sum_{s=1}^T\E_{z\sim\calD_s}[\ell(w,z)].
\]
它表示“如果所有任务都已经给出，并且可以联合训练，最好的模型是什么”。

如果只看已经到来的任务，也可以定义一个随时间变化的在线最优：
\[
  w_{\mathrm{glob},t}^\star
  \in
  \argmin_w
  \sum_{s=1}^t\E_{z\sim\calD_s}[\ell(w,z)].
\]
这个 \(w_{\mathrm{glob},t}^\star\) 一般会随 \(t\) 改变，因为新任务加入后，优化目标变了。

所以要区分三类“最优点”：
\begin{itemize}
  \item \(w_{\mathrm{glob},T}^\star\)：最终所有任务的全局最优，主要用于评价。
  \item \(w_{\mathrm{glob},t}^\star\)：截至第 \(t\) 个任务的全局最优，会随任务到来变化。
  \item \(w_t^\star\)：只针对第 \(t\) 个任务本身的最优点，
  \[
    w_t^\star\in\argmin_w\calL(w,Z_t).
  \]
  它不是累计任务的最优，而是单任务最优。
\end{itemize}

在论文的线性回归理论部分，作者额外假设存在一个共享真实参数 \(w^\star\)，所有任务都由
\[
  Y_t=X_tw^\star+\varepsilon_t
\]
生成。这个 \(w^\star\) 默认不随 \(t\) 改变。一般 CL 情形里不同任务可以有不同 \(w_t^\star\)，但线性理论为了得到闭式收敛界，使用了共享 \(w^\star\) 的简化模型。

\subsection{为什么最优解可能不唯一，最小范数是什么意思}

最优解不唯一的意思是：可能有很多个 \(w\) 都让损失达到同样的最小值。
最简单的例子是欠定线性回归。假设只有一个方程
\[
  w_1+w_2=1.
\]
那么
\[
  (w_1,w_2)=(1,0),(0,1),(1/2,1/2)
\]
都能完全拟合数据。它们都是最优解。

范数就是“大小”。例如向量 \(w=(w_1,w_2)\) 的二范数是
\[
  \norm{w}_2=\sqrt{w_1^2+w_2^2}.
\]
如果最优解有很多个，最小范数解就是在所有最优解里选长度最短的那个：
\[
  \argmin_w \norm{w}_2
  \quad
  \text{s.t. } w\text{ 已经达到最小损失}.
\]
这样做可以让目标模型唯一化，也常常对应更平滑或更简单的模型。

\subsection{论文到底用不用平方损失}

论文前面的系统模型写得比较一般，用任意损失 \(\ell(w,z)\) 表达 CL 和攻击框架。后面的严格证明，尤其是 T2T 和鲁棒特征防御的闭式推导，使用的是平方损失：
\[
  \ell(w,(x,y))=\frac12\norm{x^\top w-y}_2^2.
\]
原因是平方损失的梯度、Hessian、线性回归递推都能写成显式矩阵公式。

\subsection{\(H_t\) 的维度怎么确定}

如果参数 \(w\in\R^p\) 是一个 \(p\) 维向量，那么
\[
  H_t\in\R^{p\times p}.
\]
这是由正则化项的形状决定的：
\[
  \norm{w-w_{t-1}}_{H_t}^2
  =
  (w-w_{t-1})^\top H_t(w-w_{t-1}).
\]
其中 \(w-w_{t-1}\) 是 \(p\times1\)，左边转置是 \(1\times p\)。为了最后得到一个标量，\(H_t\) 必须是 \(p\times p\)。

在线性回归多输出情形，原始参数 \(w\in\R^{p_1\times p_2}\) 是矩阵。向量化后
\[
  \vect(w)\in\R^{p_1p_2},
\]
所以这时
\[
  H_t\in\R^{p_1p_2\times p_1p_2}.
\]

\subsection{全局超额风险和正则化更新公式的关系}

全局超额风险是\textbf{评价指标}：
\[
  \calR_T(w_T)
  =
  \text{模型 }w_T\text{ 比最终全局最优差多少}.
\]
正则化更新公式是\textbf{训练算法}：
\[
  w_t\in\argmin_w
  \left\{
  \calL(w,Z_t)+\frac12\norm{w-w_{t-1}}_{H_t}^2
  \right\}.
\]
二者不是同一个东西。理想上我们想直接最小化全局风险，但 CL 中旧任务数据和未来任务不可用，所以算法只能用“当前任务经验损失 + 不要偏离旧模型太远”的方式间接控制最终全局风险。

\subsection{label poisoning 是训练阶段攻击吗}

是的，label poisoning 是训练阶段攻击。攻击者不是在测试时改输入，而是在模型训练时提供被污染的数据或标签。

现实例子包括：
\begin{itemize}
  \item 众包标注平台中，恶意标注者提交错误标签。
  \item 联邦学习或在线学习中，某些客户端上传带错误标签的数据。
  \item 用户反馈系统中，模型把用户提供的反馈当作新训练标签，而攻击者伪造反馈。
  \item 医疗、BCI、传感器系统中，训练数据标签来自人工或外部系统，攻击者污染标签记录。
\end{itemize}
所以“正常使用的模型是否接受用户 label”取决于应用。有些系统只接收人工标注数据，有些系统会把用户反馈、交互结果或弱监督信号纳入后续训练。一旦训练管线信任这些标签，就存在 label poisoning 风险。

\subsection{零均值扰动是否代表分布不变}

不代表。零均值只表示平均方向不变：
\[
  \E[\tilde\eta_t\mid\calI_t]=0.
\]
它不保证整个分布不变。一个简单例子：
\[
  \tilde\eta=
  \begin{cases}
  +10,&\text{概率 }1/2,\\
  -10,&\text{概率 }1/2.
  \end{cases}
\]
它的均值是 0，但方差很大，数据会变得更分散、更噪。对训练来说，方差增大会让梯度更抖、估计更不稳定、收敛更慢，甚至让模型沿脆弱方向产生大误差。

所以论文的 non-shifted attack 不是说数据分布完全不变，而是说“不改变条件均值”。它仍然可以改变方差、尾部、相关结构等。

\subsection{攻击者知道什么}

论文采用强攻击者假设：攻击者知道当前 CL 模型状态和当前数据，可以策略性选择扰动。记号
\[
  \calI_t
\]
表示攻击者选择扰动时可用的信息，通常包括历史、当前模型状态和当前任务数据。

这不一定表示模型系统真的保存了所有旧数据；它是理论中的强对手建模。用强对手证明防御有效，结论更保守。如果现实攻击者不知道这么多，防御通常只会更容易。

\subsection{这些矩阵都是 Hessian 吗}

不是。论文里有几类矩阵：
\begin{itemize}
  \item \(Q_t=\nabla_{ww}^2\calL(w_t^\star,Z_t)\)：这是 Hessian，表示任务损失对参数 \(w\) 的二阶导。
  \item \(H_t\)：这是正则化矩阵，不是从数据损失直接求 Hessian 得到的，但它是正则化项的 Hessian。
  \item \(S_t=Q_t+H_t\)：这是总目标的局部 Hessian。
  \item \(G_t=\nabla_{wZ}^2\calL(w_t^\star,Z_t)\)：这是混合二阶导，表示数据变化如何影响参数梯度，不是单纯对 \(w\) 的 Hessian。
  \item \(P_t=\sum_{s=1}^tQ_s\)：这是累计 Hessian，用来衡量模型偏移对所有已见任务损失的总影响。
\end{itemize}

\subsection{KKT 条件是什么}

KKT 条件是带约束优化的一阶最优性条件。它把“目标函数最小”和“约束满足”同时写出来。
例如问题
\[
  \min_x f(x)\quad\text{s.t.}\quad g(x)\le0
\]
的拉格朗日函数是
\[
  \mathcal J(x,\lambda)=f(x)+\lambda g(x),\qquad \lambda\ge0.
\]
KKT 条件包括：
\begin{itemize}
  \item stationarity：\(\nabla_x\mathcal J(x,\lambda)=0\)；
  \item primal feasibility：\(g(x)\le0\)；
  \item dual feasibility：\(\lambda\ge0\)；
  \item complementary slackness：\(\lambda g(x)=0\)。
\end{itemize}
Proposition 5.2 中的 protected set 就是通过 KKT 条件判断哪些方向的攻击敏感度约束是“活跃”的。

\subsection{为什么向量 Taylor 展开要有转置}

若 \(w\in\R^p\)，则 \(w-w_0\) 是 \(p\times1\) 向量，Hessian \(\nabla^2f(w_0)\) 是 \(p\times p\) 矩阵。二阶项必须是标量，所以写成
\[
  (w-w_0)^\top\nabla^2f(w_0)(w-w_0).
\]
维度是
\[
  (1\times p)(p\times p)(p\times1)=1\times1.
\]
这就是为什么前面要有转置。

\subsection{一阶最优性等式不能随便写}

等式
\[
  \nabla_w\calL(w_t,Z_t+\eta_t)+H_t(w_t-w_{t-1})=0
\]
不是随便写的，也不是因为“各阶导数都为 0”。它只表示：\(w_t\) 是下面这个正则化训练目标的极小点：
\[
  J_t(w)
  =
  \calL(w,Z_t+\eta_t)
  +\frac12(w-w_{t-1})^\top H_t(w-w_{t-1}).
\]
为什么会突然出现这个目标？因为这就是论文一开始定义的正则化式 CL 更新：
\[
  \text{当前任务损失}+\text{不要离上一轮模型太远的惩罚}.
\]
没有攻击时当前数据是 \(Z_t\)；有投毒时训练系统实际收到的是 \(Z_t+\eta_t\)，所以当前任务损失从 \(\calL(w,Z_t)\) 变成 \(\calL(w,Z_t+\eta_t)\)。正则化项不变，因为它是算法自己加的稳定项。

对 \(J_t(w)\) 求梯度：
\[
  \nabla_wJ_t(w)
  =
  \nabla_w\calL(w,Z_t+\eta_t)
  +
  H_t(w-w_{t-1}).
\]
这个求导分两部分：
\begin{itemize}
  \item 第一部分 \(\calL(w,Z_t+\eta_t)\) 对 \(w\) 的导数，就是 \(\nabla_w\calL(w,Z_t+\eta_t)\)。
  \item 第二部分是二次惩罚。令 \(u=w-w_{t-1}\)，则
  \[
    \frac12u^\top H_tu
  \]
  对 \(u\) 的导数是 \(H_tu\)；又因为 \(u=w-w_{t-1}\)，对 \(w\) 求导还是 \(H_t(w-w_{t-1})\)。
\end{itemize}
在最优点 \(w=w_t\)，一阶最优性条件给出
\[
  \nabla_wJ_t(w_t)=0,
\]
于是得到论文的等式。这里为 0 的是\textbf{一阶导数/梯度}，不是所有高阶导数。二阶导数通常不为 0，它反而给出曲率信息。

第二项的 Hessian 为什么是 \(H_t\)？因为
\[
  \nabla_w\frac12(w-w_{t-1})^\top H_t(w-w_{t-1})
  =
  H_t(w-w_{t-1}),
\]
再对 \(w\) 求一次导数就是
\[
  \nabla_w^2\frac12(w-w_{t-1})^\top H_t(w-w_{t-1})
  =
  H_t.
\]

\subsection{\(G_t\) 的作用用一个线性回归例子看}

平方损失下
\[
  \calL(w,X,Y)=\frac1{2n}\norm{Xw-Y}_2^2.
\]
对 \(w\) 求梯度：
\[
  \nabla_w\calL(w,X,Y)=\frac1nX^\top(Xw-Y).
\]
如果 label 被投毒，\(Y\) 变成 \(Y+\eta\)，那么
\[
  \nabla_w\calL(w,X,Y+\eta)
  =
  \frac1nX^\top(Xw-Y-\eta)
  =
  \nabla_w\calL(w,X,Y)-\frac1nX^\top\eta.
\]
所以 label 扰动对梯度的影响是
\[
  -\frac1nX^\top\eta.
\]
这就是 \(G_t\vect(\eta_t)\) 在一般公式里的作用：它把“数据扰动”转换成“梯度扰动”。

\subsection{奇异、不满秩会造成什么影响}

矩阵奇异表示它不可逆。对方阵 \(A\)，奇异等价于
\[
  \det(A)=0.
\]
\(\det(A)\) 叫行列式，是一个把方阵映射成数字的运算。二维矩阵
\[
  A=
  \begin{bmatrix}
  a&b\\
  c&d
  \end{bmatrix}
\]
的行列式是
\[
  \det(A)=ad-bc.
\]
几何上，\(|\det(A)|\) 表示这个矩阵把面积放大了多少倍。如果 \(\det(A)=0\)，说明某个方向被压扁，面积变成 0，信息丢了，所以矩阵没有逆。

奇异也等价于存在非零向量 \(v\) 使得
\[
  Av=0.
\]
这说明某些方向被矩阵压成 0，信息丢失。

不满秩表示矩阵的独立行或列不够多。比如
\[
  A=
  \begin{bmatrix}
  1&2\\
  2&4
  \end{bmatrix}
\]
第二行是第一行的 2 倍，所以秩只有 1，不满秩，也不可逆。

在优化里，Hessian 奇异意味着某些方向曲率为 0，损失在这些方向很平坦。结果是：
\begin{itemize}
  \item 最优解可能不唯一；
  \item 普通逆矩阵不存在；
  \item 参数扰动方程不能用 \(S_t^{-1}\) 解，只能用伪逆 \(S_t^+\)。
\end{itemize}

\subsection{为什么把 \(w\) 拉直成 \(\vect(w)\)}

多输出线性回归中，\(w\) 本来是矩阵：
\[
  w\in\R^{p_1\times p_2}.
\]
当然可以直接用矩阵微积分推导，但写均值、协方差、Hessian、minimax 时会很繁琐。把矩阵拉直成向量
\[
  \vect(w)\in\R^{p_1p_2}
\]
后，就可以统一使用普通向量公式：
\[
  \Cov(\vect(w_t-w^\star)),\qquad
  H_t\in\R^{p_1p_2\times p_1p_2}.
\]
这只是记号和计算上的方便，不是改变模型本身。

\subsection{为什么要做奇异值分解}

奇异值分解把矩阵 \(A\) 写成
\[
  A=U\Sigma V^\top.
\]
可以理解成三步：
\[
  \text{先旋转输入 }(V^\top),\quad
  \text{按不同轴缩放 }(\Sigma),\quad
  \text{再旋转输出 }(U).
\]
如果
\[
  \Sigma=\diag(5,1),
\]
那么第一个奇异方向会被放大 5 倍，第二个方向只放大 1 倍。攻击者要最大化 \(\norm{A\eta}\)，当然会把预算放到第一个方向。

这就是论文说“最大奇异方向”的原因。它不是为了分解矩阵好看，而是为了解
\[
  \max_{\E\norm{\eta}^2\le M}\E\norm{A\eta}^2.
\]

\subsection{噪声协方差公式是什么意思}

论文假设
\[
  \E[\varepsilon_t^i(\varepsilon_t^i)^\top\mid x_t^i]=\sigma^2I.
\]
这里 \(\varepsilon_t^i\) 是第 \(i\) 个样本的输出噪声向量。如果输出是一维，公式退化成
\[
  \E[(\varepsilon_t^i)^2\mid x_t^i]=\sigma^2.
\]
如果输出是多维，它表示噪声协方差是 \(\sigma^2I\)，即：
\begin{itemize}
  \item 每个输出维度的噪声方差都是 \(\sigma^2\)；
  \item 不同输出维度之间不相关；
  \item 噪声强度不依赖具体输入 \(x_t^i\)。
\end{itemize}

\subsection{为什么写 \(P_t\)}

在 Proposition 2.4 中，
\[
  P_t=\sum_{\tau=1}^tQ_\tau,
  \qquad
  Q_\tau=\nabla_{ww}^2\calL(w_\tau^\star,Z_\tau).
\]
它来自把所有已见任务的经验损失都做二阶 Taylor 展开：
\[
  \sum_{\tau=1}^t
  \frac12(w_t-w_\tau^\star)^\top Q_\tau(w_t-w_\tau^\star).
\]
如果只看攻击造成的共同参数偏移 \(\delta\)，这些二次项会累加成
\[
  \frac12\delta^\top
  \left(\sum_{\tau=1}^tQ_\tau\right)
  \delta
  =
  \frac12\norm{\delta}_{P_t}^2.
\]
所以 \(P_t\) 是“所有旧任务对参数偏移的累计惩罚权重”。

\subsection{从正则化目标到论文式(4)的思路}

论文式(4)是攻防问题：
\[
  \min_{H_t}\max_{\hat\eta_t,\tilde\eta_t}
  \calR_t(w_t)
  \quad
  \text{s.t. }
  \E[\norm{\hat\eta_t+\tilde\eta_t}_F^2]\le M.
\]
它不是从正则化目标“直接求导”得到的，而是把两件事合起来建模：
\begin{itemize}
  \item 给定 \(H_t\) 和攻击扰动，模型 \(w_t\) 由正则化训练目标产生；
  \item 防御者想让最终风险小，攻击者想让最终风险大。
\end{itemize}
所以正则化目标决定 \(w_t\) 如何依赖 \(H_t,\eta_t\)；式(4)再把这种依赖放进风险 \(\calR_t(w_t)\)，形成外层的 minimax 游戏。

\subsection{攻击预算为什么用 Frobenius 范数}

一个任务里有 \(n_t\) 个样本，每个样本的标签可能是多维的，所以 label perturbation \(\eta_t\) 通常是矩阵：
\[
  \eta_t\in\R^{n_t\times p_2}.
\]
Frobenius 范数
\[
  \norm{\eta_t}_F^2=\sum_{i,j}(\eta_t^{i,j})^2
\]
就是把所有样本、所有输出维度上的扰动能量加起来。它自然表示“这次攻击总共用了多少扰动预算”。

\subsection{攻击预算为什么会放到最大奇异方向}

把 Proposition 2.4 的主导攻击项抽象成
\[
  \max_{\eta}\;\E\norm{A\eta}_2^2
  \quad
  \text{s.t.}\quad
  \E\norm{\eta}_2^2\le M,
\]
其中
\[
  A=P_t^{1/2}S_t^+G_t.
\]
若攻击是随机零均值扰动，令协方差
\[
  \Sigma=\E[\eta\eta^\top],
  \qquad
  \tr(\Sigma)=\E\norm{\eta}_2^2\le M.
\]
目标函数变成
\[
  \E\norm{A\eta}_2^2
  =
  \E[\eta^\top A^\top A\eta]
  =
  \tr(A^\top A\Sigma).
\]
设 \(A^\top A\) 的最大特征值是 \(\lambda_{\max}\)，对应单位特征向量是 \(v_1\)。因为 \(\Sigma\succeq0\)，在预算 \(\tr(\Sigma)\le M\) 下，最大化
\[
  \tr(A^\top A\Sigma)
\]
的最优选择是
\[
  \Sigma=Mv_1v_1^\top.
\]
也就是说，攻击者把全部方差预算放在 \(v_1\) 方向。这个 \(v_1\) 正是 \(A\) 的最大右奇异方向，因为 \(A^\top A\) 的特征向量就是 \(A\) 的右奇异向量。

\section{命题 2.4 / 附录命题 A.1：攻击目标的近似公式}

\subsection{结论形式}

令
\[
  \calL(w,Z_t)=\frac{1}{n_t}\sum_{z\in D_t}\ell(w,z),
  \qquad
  w_t^\star\in\argmin_w\calL(w,Z_t),
\]
并记
\[
  P_t=\sum_{s=1}^t\nabla_{ww}^2\calL(w_s^\star,Z_s).
\]
对 non-shifted 攻击 \(\tilde\eta_t\)，论文得到近似目标：
\begin{align}
  \max_{\tilde\eta_t}\quad
  &\frac12
  \E\!\left[
  \left\|
  (\nabla_{ww}^2\calL_t+H_t)^+
  \nabla_{wZ}^2\calL_t
  \vect(\tilde\eta_t)
  \right\|_{P_t}^2
  \right]
  +\text{高阶余项}
  \notag\\
  \st&
  \E[\fnorm{\tilde\eta_t}^2]\le M .
  \label{eq:attack-surrogate}
\end{align}
其中 \(\calL_t\) 是 \(\calL(w_t^\star,Z_t)\) 的简写。

这条公式表达了一个核心事实：攻击者会把预算放到
\[
  P_t^{1/2}
  (\nabla_{ww}^2\calL_t+H_t)^+
  \nabla_{wZ}^2\calL_t
\]
的最大奇异方向上，因为这个方向让相同扰动造成最大的累计风险增量。

\subsection{证明拆解}

\paragraph{第一步：用经验风险替代真实风险。}
真实风险 \(\calR_t(w_t)\) 依赖未知分布 \(\calD_1,\dots,\calD_t\)。CL 中旧数据也不可访问。因此论文先用已经出现任务的经验损失和来近似：
\[
  \sum_{\tau=1}^t\calL(w_t,Z_\tau).
\]

\paragraph{第二步：在每个任务最优点 \(w_\tau^\star\) 附近 Taylor 展开。}
因为 \(w_\tau^\star\) 是任务 \(\tau\) 的经验损失极小点，
\[
  \nabla_w\calL(w_\tau^\star,Z_\tau)=0 .
\]
所以二阶展开给出
\begin{align}
  \sum_{\tau=1}^t\calL(w_t,Z_\tau)
  =
  \sum_{\tau=1}^t\calL(w_\tau^\star,Z_\tau)
  +\frac12\sum_{\tau=1}^t
  (w_t-w_\tau^\star)^\top Q_\tau(w_t-w_\tau^\star)
  +O\!\left(\sum_{\tau=1}^t\norm{w_t-w_\tau^\star}_2^3\right),
  \label{eq:taylor-risk}
\end{align}
其中 \(Q_\tau=\nabla_{ww}^2\calL(w_\tau^\star,Z_\tau)\)。

\paragraph{第三步：把投毒导致的更新写成局部线性方程。}
带投毒数据 \(Z_t+\eta_t\) 的一阶最优性条件是
\[
  \nabla_w\calL(w_t,Z_t+\eta_t)+H_t(w_t-w_{t-1})=0 .
\]
在 \((w_t^\star,Z_t)\) 附近展开：
\[
  \nabla_w\calL(w_t,Z_t+\eta_t)
  \approx
  \underbrace{\nabla_w\calL(w_t^\star,Z_t)}_{=0}
  +
  \underbrace{\nabla_{ww}^2\calL(w_t^\star,Z_t)}_{Q_t}
  (w_t-w_t^\star)
  +
  \underbrace{\nabla_{wZ}^2\calL(w_t^\star,Z_t)}_{G_t}
  \vect(\eta_t).
\]
第一项为 0，是因为 \(w_t^\star\) 是干净任务 \(t\) 的经验损失最小点。把这个展开代回一阶最优性条件，就得到
\begin{equation}
  Q_t(w_t-w_t^\star)
  +G_t\vect(\eta_t)
  +H_t(w_t-w_{t-1})
  +O(\norm{w_t-w_t^\star}_2^2+\fnorm{\eta_t}^2)=0,
  \label{eq:linearized-foc}
\end{equation}
其中 \(G_t=\nabla_{wZ}^2\calL(w_t^\star,Z_t)\)。
下面把 \eqref{eq:linearized-foc} 整理成论文使用的形式。目标是求 \(w_t-w_\tau^\star\)，其中 \(\tau\le t\) 是任意一个历史任务。令
\[
  a_\tau:=w_t-w_\tau^\star.
\]
那么
\[
  w_t-w_t^\star
  =
  a_\tau-(w_t^\star-w_\tau^\star),
  \qquad
  w_t-w_{t-1}
  =
  a_\tau-(w_{t-1}-w_\tau^\star).
\]
代入 \eqref{eq:linearized-foc}：
\begin{align*}
  &Q_t\bigl[a_\tau-(w_t^\star-w_\tau^\star)\bigr]
  +G_t\vect(\eta_t)
  +H_t\bigl[a_\tau-(w_{t-1}-w_\tau^\star)\bigr]
  +\text{高阶项}
  =0.
\end{align*}
把含 \(a_\tau\) 的项放左边，其余项放右边：
\[
  (Q_t+H_t)a_\tau
  =
  Q_t(w_t^\star-w_\tau^\star)
  +H_t(w_{t-1}-w_\tau^\star)
  -G_t\vect(\eta_t)
  +\text{高阶项}.
\]
记 \(S_t=Q_t+H_t\)。若 \(S_t\) 可逆，就左乘 \(S_t^{-1}\)；若不可逆，就用伪逆 \(S_t^+\)。于是
\begin{align}
  w_t-w_\tau^\star
  =&
  S_t^+
  \left[
    Q_t(w_t^\star-w_\tau^\star)
    +H_t(w_{t-1}-w_\tau^\star)
    -G_t\vect(\eta_t)
  \right]
  +O(\norm{w_t-w_t^\star}_2^2+\fnorm{\eta_t}^2),
  \label{eq:update-shift}
\end{align}
这就是局部线性方程的来源。

\paragraph{第四步：代回二次风险。}
把 \eqref{eq:update-shift} 代入 \eqref{eq:taylor-risk}，所有和 \(\eta_t\) 无关的项对攻击者不重要，可以丢掉。留下的攻击相关部分是
\begin{align}
  &
  \frac12\norm{S_t^+G_t\vect(\eta_t)}_{P_t}^2
  -
  \vect(\eta_t)^\top G_t^\top S_t^+
  \sum_{\tau=1}^t
  Q_\tau S_t^+
  \left[
    Q_t(w_t^\star-w_\tau^\star)
    +H_t(w_{t-1}-w_\tau^\star)
  \right]
  \notag\\
  &\quad
  +O(\fnorm{\eta_t}^2+\norm{w_t-w_t^\star}_2^2)
  +O\!\left(\sum_{\tau=1}^t\norm{w_t-w_\tau^\star}_2^3\right).
  \label{eq:attack-expansion}
\end{align}

这一行可以按下面的方式一步步看。先定义
\[
  b_\tau
  =
  S_t^+
  \left[
    Q_t(w_t^\star-w_\tau^\star)
    +H_t(w_{t-1}-w_\tau^\star)
  \right],
  \qquad
  r_t=S_t^+G_t\vect(\eta_t).
\]
那么忽略高阶项时
\[
  w_t-w_\tau^\star=b_\tau-r_t.
\]
二次型记号是
\[
  \norm{x}_{Q_\tau}^2=x^\top Q_\tau x.
\]
所以 Taylor 二次项变成
\[
  \frac12\sum_{\tau=1}^t\norm{b_\tau-r_t}_{Q_\tau}^2.
\]
展开一个平方：
\[
  \norm{b_\tau-r_t}_{Q_\tau}^2
  =
  (b_\tau-r_t)^\top Q_\tau(b_\tau-r_t)
  =
  b_\tau^\top Q_\tau b_\tau
  -2r_t^\top Q_\tau b_\tau
  +r_t^\top Q_\tau r_t.
\]
三类项的含义是：
\begin{itemize}
  \item \(b_\tau^\top Q_\tau b_\tau\)：不含 \(\eta_t\)，攻击者无法通过当前扰动改变，分析攻击目标时可以丢掉。
  \item \(-2r_t^\top Q_\tau b_\tau\)：含一个 \(r_t\)，所以是关于 \(\eta_t\) 的线性项。
  \item \(r_t^\top Q_\tau r_t\)：含两个 \(r_t\)，所以是关于 \(\eta_t\) 的二次项。
\end{itemize}
把二次项对 \(\tau\) 求和：
\[
  \frac12\sum_{\tau=1}^t r_t^\top Q_\tau r_t
  =
  \frac12 r_t^\top
  \left(\sum_{\tau=1}^tQ_\tau\right)
  r_t
  =
  \frac12\norm{r_t}_{P_t}^2.
\]
再代回 \(r_t=S_t^+G_t\vect(\eta_t)\)，就得到
\[
  \frac12\norm{S_t^+G_t\vect(\eta_t)}_{P_t}^2.
\]
线性项对 \(\tau\) 求和后就是 \eqref{eq:attack-expansion} 中的第二项。由于 \(S_t\) 是对称矩阵，伪逆 \(S_t^+\) 也可按对称矩阵处理，所以转置记号可以简化成论文里的形式。

\paragraph{第五步：为什么 non-shifted 攻击没有线性项。}
若
\[
  \eta_t=\hat\eta_t+\tilde\eta_t,\qquad
  \E[\tilde\eta_t\mid Z_{1:t},w_{t-1}]=0,
\]
把 \(\eta_t=\tilde\eta_t\) 代入 \eqref{eq:attack-expansion}。其中第一项是二次型：
\[
  \frac12\norm{S_t^+G_t\vect(\tilde\eta_t)}_{P_t}^2.
\]
第二项是线性项，可以写成
\[
  -\vect(\tilde\eta_t)^\top a_t,
\]
其中 \(a_t\) 是由 \(G_t,S_t,Q_\tau,w_t^\star,w_\tau^\star,w_{t-1}\) 组成的、在条件信息 \(Z_{1:t},w_{t-1}\) 下已经确定的向量。于是
\[
  \E[-\vect(\tilde\eta_t)^\top a_t\mid Z_{1:t},w_{t-1}]
  =
  -\E[\vect(\tilde\eta_t)\mid Z_{1:t},w_{t-1}]^\top a_t
  =
  0.
\]
所以对 non-shifted 攻击，取期望后线性项消失，主导项只剩
\[
  \frac12
  \E\left[
  \norm{S_t^+G_t\vect(\tilde\eta_t)}_{P_t}^2
  \right],
\]
这正是 \eqref{eq:attack-surrogate}。
shifted 攻击 \(\hat\eta_t\) 会保留线性项，所以它更危险：它能直接把模型均值推向错误方向。

\section{定理 3.1：为什么有些攻击根本防不住}

\subsection{定理含义}

论文证明：存在一些问题实例，使得任意正则化式 CL 方法 \eqref{eq:regularized-cl} 都无法保证收敛。两类攻击是：
\begin{itemize}
  \item frequent + shifted；
  \item frequent + unbounded + non-shifted。
\end{itemize}
证明只需要构造一维线性回归反例。

反例证明的原理是：如果一个结论声称“任何正则化式 CL 方法都能防住”，那么只要找出一个非常简单的问题，使得任意 \(H_t\) 选择都会失败，就能推翻这个声称。一维问题已经足够，因为正则化更新会退化成一个简单的加权平均：
\[
  \text{新模型}
  =
  \text{当前数据权重}\times\text{当前观测}
  +
  \text{旧模型权重}\times\text{旧模型}.
\]
攻击者正是利用这个加权平均：如果算法不断吸收新数据，就会吸收攻击；如果算法不吸收新数据，又无法消除初始误差。

\subsection{一维反例}

设
\[
  p_1=p_2=n_t=1,\qquad x_t=1,\qquad y_t=w^\star+\varepsilon_t,
\]
平方损失 \(\ell(w,z_t)=(w-y_t)^2\)。
这里 \(w^\star\) 是真实的一维参数。因为 \(x_t=1\)，干净数据满足 \(y_t=w^\star+\varepsilon_t\)，所以如果没有噪声和攻击，最好的模型就是 \(w=w^\star\)。

若任务 \(t\) 被投毒，训练标签变成 \(y_t+\eta_t\)。正则化更新目标是
\[
  \min_w
  \left\{
  (w-(y_t+\eta_t))^2
  +\frac12H_t(w-w_{t-1})^2
  \right\}.
\]
对 \(w\) 求导并令其为 0：
\[
  2(w-(y_t+\eta_t))+H_t(w-w_{t-1})=0.
\]
展开：
\[
  (2+H_t)w=2(y_t+\eta_t)+H_tw_{t-1}.
\]
因此
\[
  w_t
  =
  \frac{2}{2+H_t}(y_t+\eta_t)
  +
  \frac{H_t}{2+H_t}w_{t-1}.
\]
论文说“吸收常数因子”，意思是把常数 \(2\) 合并进正则化参数的记号里。因为这里要证明的是不可能性，常数倍不影响结论，于是等价写成
\[
  w_t
  =
  \frac{1}{1+H_t}(y_t+\eta_t)
  +
  \frac{H_t}{1+H_t}w_{t-1}.
\]
定义
\[
  c_t=\frac{H_t}{1+H_t}\in[0,1],
\]
则 \(1-c_t=1/(1+H_t)\)。把 \(y_t=w^\star+\varepsilon_t\) 代入：
\[
  w_t
  =
  (1-c_t)(w^\star+\varepsilon_t+\eta_t)+c_tw_{t-1}.
\]
两边减去 \(w^\star\)：
\[
  w_t-w^\star
  =
  c_t(w_{t-1}-w^\star)
  +(1-c_t)(\varepsilon_t+\eta_t)
  +(1-c_t)w^\star+c_tw^\star-w^\star.
\]
最后三项中关于 \(w^\star\) 的部分抵消，因为
\[
  (1-c_t)w^\star+c_tw^\star-w^\star=0.
\]
因此
\begin{equation}
  w_t-w^\star
  =
  c_t(w_{t-1}-w^\star)
  +(1-c_t)(\varepsilon_t+\eta_t).
  \label{eq:scalar-recursion}
\end{equation}
这个递推是整条证明的核心。

\subsection{frequent shifted 攻击}

令攻击每个任务，并取常数偏移 \(\eta_t=\Delta\neq 0\)。
设
\[
  m_t=\E[w_t-w^\star],
\]
由 \eqref{eq:scalar-recursion} 得
\[
  m_t=c_tm_{t-1}+(1-c_t)\Delta .
\]
减去 \(\Delta\)：
\[
  m_t-\Delta=c_t(m_{t-1}-\Delta),
\]
这个式子每往前展开一次，就多乘一个 \(c\)。例如：
\[
  m_T-\Delta
  =
  c_T(m_{T-1}-\Delta)
  =
  c_Tc_{T-1}(m_{T-2}-\Delta)
  =
  \cdots
  =
  \left(\prod_{t=1}^Tc_t\right)(m_0-\Delta).
\]
而 \(m_0=\E[w_0-w^\star]=w_0-w^\star\)，所以
\[
  m_T=\Delta+\left(\prod_{t=1}^Tc_t\right)(w_0-w^\star-\Delta).
\]
现在分两种情况：
\begin{itemize}
  \item 若 \(\prod_{t=1}^Tc_t\not\to0\)，则即使无攻击无噪声，初始偏差也不能被完全遗忘，算法本来就不保证收敛。
  \item 若 \(\prod_{t=1}^Tc_t\to0\)，则上式第二项趋于 0，所以在 shifted 攻击下 \(m_T\to\Delta\)，于是
  \[
    \E[(w_T-w^\star)^2]\ge (\E[w_T-w^\star])^2\to\Delta^2>0 .
  \]
  这里用的是
  \[
    \E[X^2]=\operatorname{Var}(X)+(\E X)^2\ge(\E X)^2.
  \]
  令 \(X=w_T-w^\star\)。如果均值已经趋向 \(\Delta\)，那么平方误差的期望至少趋向 \(\Delta^2\)，不可能收敛到 0。
\end{itemize}
因此无论正则化策略如何选，总有一种情形无法收敛。

\subsection{frequent unbounded non-shifted 攻击}

现在攻击仍然每个任务，但取零均值独立扰动：
\[
  \E[\tilde\eta_t]=0,\qquad
  \E[\tilde\eta_t^2]=M_T,\qquad
  M_T=\Omega(T).
\]
\(\E[\tilde\eta_t^2]=M_T\) 在一维零均值情况下就是“攻击扰动的方差等于 \(M_T\)”。\(M_T=\Omega(T)\) 的意思是：存在常数 \(c>0\) 和足够大的 \(T\)，使得
\[
  M_T\ge cT.
\]
也就是说，单次攻击预算至少和任务总数 \(T\) 同阶增长，所以这是 unbounded attack。

展开 \eqref{eq:scalar-recursion}：
\[
  w_T-w^\star
  =
  \left(\prod_{r=1}^Tc_r\right)(w_0-w^\star)
  +
  \sum_{s=1}^T
  \beta_{s,T}(\varepsilon_s+\tilde\eta_s),
\]
其中
\[
  \beta_{s,T}=(1-c_s)\prod_{r=s+1}^Tc_r .
\]
\(\beta_{s,T}\) 来自反复代入递推。比如写到 \(T=3\)：
\begin{align*}
  w_3-w^\star
  =&c_3(w_2-w^\star)+(1-c_3)(\varepsilon_3+\tilde\eta_3)\\
  =&c_3c_2(w_1-w^\star)
    +c_3(1-c_2)(\varepsilon_2+\tilde\eta_2)
    +(1-c_3)(\varepsilon_3+\tilde\eta_3)\\
  =&c_3c_2c_1(w_0-w^\star)
    +c_3c_2(1-c_1)(\varepsilon_1+\tilde\eta_1)\\
  &\quad
    +c_3(1-c_2)(\varepsilon_2+\tilde\eta_2)
    +(1-c_3)(\varepsilon_3+\tilde\eta_3).
\end{align*}
因此第 \(s\) 步噪声前面的系数就是“进入时的 \(1-c_s\)”乘上“之后每一步保留下来的 \(c\)”。

\(\beta_{s,T}\) 的含义是：第 \(s\) 步噪声进入模型时先乘当前数据权重 \(1-c_s\)，之后从 \(s+1\) 到 \(T\) 每一步又被旧模型保留权重 \(c_{s+1},\dots,c_T\) 继续传递，所以总权重是
\[
  (1-c_s)c_{s+1}c_{s+2}\cdots c_T.
\]
当 \(s=T\) 时，后面已经没有 \(c_{T+1},\dots,c_T\) 可以乘。数学上把这种“没有任何项可乘”的乘积定义为 1，叫空乘积：
\[
  \prod_{r=T+1}^Tc_r=1.
\]
所以
\[
  \beta_{T,T}=1-c_T.
\]
这些权重满足望远镜恒等式
\[
  \sum_{s=1}^T\beta_{s,T}
  =
  1-\prod_{r=1}^Tc_r .
\]
为什么叫望远镜？因为
\[
  \beta_{s,T}
  =
  (1-c_s)\prod_{r=s+1}^Tc_r
  =
  \prod_{r=s+1}^Tc_r-\prod_{r=s}^Tc_r.
\]
于是求和时中间项会一正一负相互抵消：
\begin{align*}
  \sum_{s=1}^T\beta_{s,T}
  &=
  (c_2\cdots c_T-c_1c_2\cdots c_T)
  +(c_3\cdots c_T-c_2c_3\cdots c_T)
  +\cdots
  +(1-c_T)\\
  &=
  1-c_1c_2\cdots c_T
  =
  1-\prod_{r=1}^Tc_r.
\end{align*}

如果 \(\prod c_r\not\to0\)，初始偏差仍然不消失；否则
\[
  \sum_{s=1}^T\beta_{s,T}\to1 .
\]
由 Cauchy 不等式
\[
  \sum_{s=1}^T\beta_{s,T}^2
  \ge
  \frac{\left(\sum_{s=1}^T\beta_{s,T}\right)^2}{T}
  =
  \frac{(1-\prod_{r=1}^Tc_r)^2}{T}.
\]
Cauchy 在这里用的是基础形式：
\[
  \left(\sum_{s=1}^T a_s\cdot 1\right)^2
  \le
  \left(\sum_{s=1}^T a_s^2\right)
  \left(\sum_{s=1}^T 1^2\right)
  =
  T\sum_{s=1}^T a_s^2.
\]
把 \(a_s=\beta_{s,T}\) 代入并移项，就得到上面的下界。

攻击方差给最终风险贡献至少
\[
  M_T\sum_{s=1}^T\beta_{s,T}^2
  \ge
  M_T\frac{(1-\prod_{r=1}^Tc_r)^2}{T}.
\]
因为 \(M_T=\Omega(T)\)，右端不会趋于 0。定理得证。

\begin{remark}
这个反例很重要。正则化 CL 必须靠 \(c_t<1\) 逐步忘掉初始偏差；但一旦每一步都混入有偏或巨大方差攻击，这种“持续吸收新数据”的机制本身就会把攻击吸收进去。
\end{remark}

\section{T2T 防御：少量但强攻击的检测与证明}

\subsection{先不看公式：T2T 到底想做什么}

T2T 的问题不是“模型有没有变”，而是“最近两个任务里有没有异常投毒”。只看
\[
  w_t-w_{t-1}
\]
不够，因为这个更新可能同时混有三种来源：
\begin{itemize}
  \item 当前任务 \(t\) 的正常学习信号；
  \item 当前或上一任务的投毒信号；
  \item 更早历史模型 \(w_{t-2}\) 留下的影响。
\end{itemize}
第三种历史影响会干扰判断。T2T 的核心就是：用连续两次更新互相对照，把共同的历史影响抵消掉，只留下最近两个任务造成的局部变化。

可以把两次更新写成一种口语化形式：
\[
  \text{当前更新}
  \approx
  A_t\times\text{旧历史状态}
  +\text{最近任务差异}
  +\text{最近投毒},
\]
\[
  \text{上一次更新}
  \approx
  B_t\times\text{旧历史状态}
  +\text{上一任务投毒}.
\]
这里的 \(A_t\) 和 \(B_t\) 不是随机构造的变量，它们是由正则化更新的线性化推出来的\textbf{历史传播矩阵}：
\begin{itemize}
  \item \(A_t\)：旧状态 \(w_{t-2}\) 经过任务 \(t-1\) 再影响任务 \(t\) 的路径。
  \item \(B_t\)：旧状态 \(w_{t-2}\) 影响上一轮更新 \(w_{t-1}-w_{t-2}\) 的路径。
\end{itemize}
换句话说，\(A_t\) 和 \(B_t\) 描述的是“同一个历史状态分别在两次更新里长成什么样”。

那为什么检测分数要乘 \(D_t^1,D_t^2\)？因为当前更新里的历史项是 \(A_t\times\text{旧历史状态}\)，上一更新里的历史项是 \(B_t\times\text{旧历史状态}\)。它们形状不同，不能直接相减。于是我们找两个线性变换 \(D_t^1,D_t^2\)，让
\[
  D_t^1A_t=D_t^2B_t.
\]
这样一来，两次更新经过 \(D_t^1,D_t^2\) 处理后，历史项变成同一个东西，相减时就消掉了。

所以 T2T 检测分数的逻辑是：
\[
  \text{检测分数}
  =
  \norm{
  \text{处理后的当前更新}
  -
  \text{处理后的上一次更新}
  }.
\]
如果没有攻击，这个差值主要来自相邻任务本身的正常差异和随机噪声，应该比较稳定；如果最近两个任务中有强投毒，这个差值会突然变大。

论文中的 \(C_t\) 也不是随机构造。它是为了在满足
\[
  D_t^1A_t=D_t^2B_t
\]
的同时，让 \(D_t^1\) 尽量像 \(A_t^+\)，让 \(D_t^2\) 尽量像 \(B_t^+\)。直观上，\(A_t^+\) 和 \(B_t^+\) 是“尽量把 \(A_t,B_t\) 的影响反解回去”的工具；但只用伪逆不一定能让历史项完全抵消，所以论文加了 \(C_t\) 来修正。

\((I-B_t)\) 的作用也不是普通意义上“把长度归一化”。这里 \(I\) 是单位矩阵，\(I-B_t\) 表示“保留原向量，再减去 \(B_t\) 能解释的那部分”。可以把它理解成一个线性过滤器：它削弱上一任务动力学已经能解释的分量，让检测分数更关注新出现的局部异常。论文把它放在 \(D_t^1,D_t^2\) 前面，是为了让两个投影后的更新处在更可比较的局部空间里。

如果你说的公式 16、17、18 对应这一段，那么它们的关系应这样读：
\begin{itemize}
  \item 检测分数公式：定义“比较什么”。
  \item \(D_t^1,D_t^2,C_t\) 公式：定义“如何把两次更新变到同一空间再比较”。
  \item \(A_t,B_t\) 公式：说明“历史状态在两次更新里分别如何传播”。
\end{itemize}
先有检测目标，再有抵消历史项的需求，最后才有这些矩阵公式。

\subsection{检测分数的设计}

少量但强攻击的关键困难是：攻击任务未知，而且 shifted 攻击可以非常大。
最自然的防御是检测异常任务并回滚。但只看
\[
  \norm{w_t-w_{t-1}}
\]
会有归因问题：大更新可能来自当前任务被攻击，也可能来自当前良性任务在纠正过去被攻击造成的偏差。

T2T 使用两个连续更新：
\[
  w_t-w_{t-1},\qquad w_{t-1}-w_{t-2}.
\]
检测分数为
\begin{equation}
  d_t=
  \norm{
  D_t^1(w_t-w_{t-1})-D_t^2(w_{t-1}-w_{t-2})
  }.
  \label{eq:dt}
\end{equation}
其中
\begin{align}
  D_t^1&=(I-B_t)(A_t^+-C_tA_t^\top),\\
  D_t^2&=(I-B_t)(B_t^++C_tB_t^\top),\\
  C_t&=(A_t^+A_t-B_t^+B_t)(A_t^\top A_t+B_t^\top B_t)^+,
  \label{eq:Dt}
\end{align}
而
\begin{align}
  A_t&=(\nabla_{ww}^2\calL_t+H_t)^{-1}
      \nabla_{ww}^2\calL_t
      (\nabla_{ww}^2\calL_{t-1}+H_{t-1})^{-1}H_{t-1},\\
  B_t&=(\nabla_{ww}^2\calL_{t-1}+H_{t-1})^{-1}
      \nabla_{ww}^2\calL_{t-1}.
\end{align}

\subsection{引理 4.1：为什么这个分数能隔离最近两个任务}

引理说：
\begin{align}
  D_t^1(w_t-w_{t-1})-D_t^2(w_{t-1}-w_{t-2})
  =
  E_t^1(w_t^\star-w_{t-1}^\star)
  +E_t^2\vect(\eta_{t-1})
  +E_t^3\vect(\eta_t)
  +\mathcal E_t,
  \label{eq:t2t-lemma}
\end{align}
其中
\[
  \mathcal E_t
  =
  O(\norm{w_t-w_t^\star}^2)
  +O(\norm{w_{t-1}-w_{t-1}^\star}^2)
  +O(\norm{\eta_t}^2+\norm{\eta_{t-1}}^2).
\]
最重要的是：右边没有 \(w_{t-2}\) 的历史项。

\subsection{证明关键：构造 \(D_t^1A_t=D_t^2B_t\)}

从线性化更新可得到
\begin{align}
  w_t-w_{t-1}
  =&
  -A_t(w_{t-2}-w_t^\star)
  +S_t^+Q_tS_{t-1}^+Q_{t-1}(w_t^\star-w_{t-1}^\star)
  \notag\\
  &+
  S_t^+Q_tS_{t-1}^+G_{t-1}\vect(\eta_{t-1})
  -S_t^+G_t\vect(\eta_t)
  +\text{高阶项},
  \label{eq:t2t-update1}\\
  w_{t-1}-w_{t-2}
  =&
  -B_t(w_{t-2}-w_{t-1}^\star)
  -S_{t-1}^+G_{t-1}\vect(\eta_{t-1})
  +\text{高阶项}.
  \label{eq:t2t-update2}
\end{align}
两式都含有历史状态 \(w_{t-2}\)。为了消掉它，需要
\[
  D_t^1A_t=D_t^2B_t .
\]
论文把 \(D_t^1,D_t^2\) 设计成下列最小二乘投影问题的解：
\[
  \min_{\bar D_t^1,\bar D_t^2}
  \norm{\bar D_t^1-A_t^+}_F^2
  +\norm{\bar D_t^2-B_t^+}_F^2,
  \qquad
  \st
  \bar D_t^1A_t=\bar D_t^2B_t .
\]
解出来就是 \eqref{eq:Dt} 中的形式，再乘上 \((I-B_t)\) 做归一化。
代回 \eqref{eq:t2t-update1}--\eqref{eq:t2t-update2} 后，历史项抵消，只剩最近两个任务的任务差异和攻击项。

\subsection{把 T2T 公式 10--25 连成一条推导链}

这一小节按“为什么需要这个公式”的顺序讲，而不是按论文出现顺序讲。

\paragraph{第一步：从单步正则化更新得到线性化更新。}
对任意任务 \(s\)，正则化训练目标是一阶最优：
\[
  \nabla_w\calL(w_s,Z_s+\eta_s)+H_s(w_s-w_{s-1})=0.
\]
在干净单任务最优点 \((w_s^\star,Z_s)\) 附近展开梯度：
\[
  \nabla_w\calL(w_s,Z_s+\eta_s)
  \approx
  Q_s(w_s-w_s^\star)+G_s\vect(\eta_s),
\]
其中
\[
  Q_s=\nabla_{ww}^2\calL(w_s^\star,Z_s),
  \qquad
  G_s=\nabla_{wZ}^2\calL(w_s^\star,Z_s),
  \qquad
  S_s=Q_s+H_s.
\]
代回最优性条件并整理：
\[
  S_s(w_s-w_a^\star)
  \approx
  Q_s(w_s^\star-w_a^\star)
  +H_s(w_{s-1}-w_a^\star)
  -G_s\vect(\eta_s).
\]
这里 \(w_a^\star\) 是我们选来作参照的某个任务最优点。左乘 \(S_s^+\) 得到
\[
  w_s-w_a^\star
  \approx
  S_s^+
  \left[
  Q_s(w_s^\star-w_a^\star)
  +H_s(w_{s-1}-w_a^\star)
  -G_s\vect(\eta_s)
  \right].
\]
这就是“线性化更新”：原本复杂的训练结果 \(w_s\)，被近似写成历史模型、任务差异、投毒扰动的线性组合。

\paragraph{第二步：为什么 \(A_t,B_t\) 是那样定义的。}
现在我们关心两次更新：
\[
  w_t-w_{t-1},
  \qquad
  w_{t-1}-w_{t-2}.
\]
先看上一轮更新。把上面的公式用于任务 \(t-1\)，并把参照点取为 \(w_{t-1}^\star\)，得到
\[
  w_{t-1}-w_{t-2}
  \approx
  -S_{t-1}^+Q_{t-1}(w_{t-2}-w_{t-1}^\star)
  -S_{t-1}^+G_{t-1}\vect(\eta_{t-1}).
\]
所以自然定义
\[
  B_t=S_{t-1}^+Q_{t-1}.
\]
\(B_t\) 的意义是：旧历史状态 \(w_{t-2}\) 如何进入上一轮更新。

再看当前更新。当前更新 \(w_t-w_{t-1}\) 里会含有 \(w_{t-1}\)，而 \(w_{t-1}\) 又由上一任务从 \(w_{t-2}\) 更新而来。把上一轮线性化结果继续代入当前轮，就会得到一条从 \(w_{t-2}\) 到 \(w_t-w_{t-1}\) 的路径：
\[
  A_t=S_t^+Q_tS_{t-1}^+H_{t-1}.
\]
\(A_t\) 的意义是：旧历史状态 \(w_{t-2}\) 经过上一轮正则化保留项 \(H_{t-1}\)，再经过当前任务曲率 \(Q_t\)，最后影响当前更新。

因此 \(A_t,B_t\) 不是随便定义的，它们分别是“同一个历史状态在当前更新和上一更新中的传播矩阵”。

\paragraph{第三步：公式 17 和 18 的来源。}
把上面的线性化代数完整代入，就得到两条关键关系：
\begin{align*}
  w_t-w_{t-1}
  \approx&
  -A_t(w_{t-2}-w_t^\star)
  +S_t^+Q_tS_{t-1}^+Q_{t-1}(w_t^\star-w_{t-1}^\star)\\
  &+
  S_t^+Q_tS_{t-1}^+G_{t-1}\vect(\eta_{t-1})
  -S_t^+G_t\vect(\eta_t),
  \\
  w_{t-1}-w_{t-2}
  \approx&
  -B_t(w_{t-2}-w_{t-1}^\star)
  -S_{t-1}^+G_{t-1}\vect(\eta_{t-1}).
\end{align*}
这两条就是 T2T 的原材料。它们告诉我们：两次更新里都带着旧历史状态 \(w_{t-2}\)，只是一个被 \(A_t\) 处理，一个被 \(B_t\) 处理。

\paragraph{第四步：为什么要 \(D_t^1A_t=D_t^2B_t\)。}
检测时要相减：
\[
  D_t^1(w_t-w_{t-1})-D_t^2(w_{t-1}-w_{t-2}).
\]
如果只看历史项，它们分别是
\[
  -D_t^1A_t(w_{t-2}-\cdots),
  \qquad
  -D_t^2B_t(w_{t-2}-\cdots).
\]
要让同一个旧历史状态在相减时抵消，必须让它们前面的线性变换一致：
\[
  D_t^1A_t=D_t^2B_t.
\]
这就是 T2T 的核心约束。没有这个约束，检测分数仍然会混入很早以前的模型状态，无法只看最近两个任务。

\paragraph{第五步：为什么变成最小二乘投影问题。}
满足 \(D_t^1A_t=D_t^2B_t\) 的 \(D_t^1,D_t^2\) 可能有很多组。论文还希望它们尽量保留原更新里的信息，而不是随便把所有东西都压成 0。自然的选择是：让 \(D_t^1\) 尽量接近 \(A_t^+\)，让 \(D_t^2\) 尽量接近 \(B_t^+\)。这就得到
\[
  \min_{\bar D_t^1,\bar D_t^2}
  \norm{\bar D_t^1-A_t^+}_F^2
  +\norm{\bar D_t^2-B_t^+}_F^2
  \quad
  \text{s.t.}\quad
  \bar D_t^1A_t=\bar D_t^2B_t.
\]
这不是为了形式复杂，而是在做一件很直观的事：既要抵消历史，又要尽量少扭曲两次更新。

\paragraph{第六步：为什么解里会出现 \(C_t\)，以及为什么一个是减号一个是加号。}
令
\[
  R_t=A_t^+A_t-B_t^+B_t.
\]
如果直接取 \(\bar D_t^1=A_t^+\)、\(\bar D_t^2=B_t^+\)，约束左边和右边相差 \(R_t\)。所以要给二者加修正量，让一个往回减，另一个往前加：
\[
  \bar D_t^1=A_t^+-R_tK_tA_t^\top,
  \qquad
  \bar D_t^2=B_t^+ +R_tK_tB_t^\top.
\]
这就是 \(D^1\) 中是减号、\(D^2\) 中是加号的原因：两个修正朝相反方向移动，才能把中间差距补上。

把它们乘上 \(A_t,B_t\)：
\[
  \bar D_t^1A_t-\bar D_t^2B_t
  =
  R_t-R_tK_t(A_t^\top A_t+B_t^\top B_t).
\]
为了让这个差为 0，需要
\[
  K_t(A_t^\top A_t+B_t^\top B_t)
  \approx I
\]
在相关子空间上成立。最小范数意义下的解就是
\[
  K_t=(A_t^\top A_t+B_t^\top B_t)^+.
\]
论文为了少写一个 \(R_tK_t\)，把两者合并记作
\[
  C_t=(A_t^+A_t-B_t^+B_t)(A_t^\top A_t+B_t^\top B_t)^+
\]
于是 \(\bar D_t^1=A_t^+-C_tA_t^\top\)，\(\bar D_t^2=B_t^++C_tB_t^\top\)。它本质上就是“用最小二乘方式补齐 \(A_t^+A_t\) 和 \(B_t^+B_t\) 的差距”。

\paragraph{第七步：代回后为什么得到 \(E_t^1,E_t^2,E_t^3\)。}
把两条线性化更新代回检测差：
\[
  D_t^1(w_t-w_{t-1})-D_t^2(w_{t-1}-w_{t-2}).
\]
历史项因为 \(D_t^1A_t=D_t^2B_t\) 消掉。剩下的项按来源分组：
\begin{itemize}
  \item 和任务差异 \(w_t^\star-w_{t-1}^\star\) 相乘的所有矩阵，合并叫 \(E_t^1\)。
  \item 和上一任务攻击 \(\vect(\eta_{t-1})\) 相乘的所有矩阵，合并叫 \(E_t^2\)。
  \item 和当前任务攻击 \(\vect(\eta_t)\) 相乘的所有矩阵，合并叫 \(E_t^3\)。
  \item Taylor 展开的高阶残差，合并叫 \(\mathcal E_t\)。
\end{itemize}
于是得到
\[
  E_t^1(w_t^\star-w_{t-1}^\star)
  +E_t^2\vect(\eta_{t-1})
  +E_t^3\vect(\eta_t)
  +\mathcal E_t.
\]
这就是你问的公式 16 的来源：不是凭空来的，而是“历史项抵消后，把剩余项按来源重新命名”。

\paragraph{第八步：为什么良性任务时 \(d_t^2=\|E_t^2e_{t-1}+E_t^3e_t\|^2\)。}
在线性回归理论分析里，良性任务没有攻击，只有自然噪声：
\[
  \mathbf n_{t-1}=0,\qquad \mathbf n_t=0.
\]
这里 \(\mathbf n_t\) 表示攻击额外加到标签上的扰动。良性时它为 0，但标签观测里仍然有自然噪声 \(\mathbf e_t\)。自然噪声和 label perturbation 一样都会通过标签通道影响梯度，所以在 T2T 分数里出现 \(\mathbf e_{t-1},\mathbf e_t\)。共享参数模型下任务最优点的系统差异项被简化，检测分数主要变成
\[
  d_t^2
  =
  \norm{E_t^2\mathbf e_{t-1}+E_t^3\mathbf e_t}_2^2.
\]
它的含义是：没有攻击时，检测分数只由最近两个任务的自然噪声造成。

\paragraph{第九步：交叉项期望为什么为 0。}
展开平方：
\begin{align*}
  d_t^2
  =&
  \norm{E_t^2\mathbf e_{t-1}}^2
  +\norm{E_t^3\mathbf e_t}^2
  +2(E_t^2\mathbf e_{t-1})^\top(E_t^3\mathbf e_t).
\end{align*}
最后一项叫交叉项。因为 \(e_{t-1}\) 和 \(e_t\) 来自不同任务，条件在特征矩阵上它们独立且均值为 0，所以
\[
  \E[\mathbf e_{t-1}\mathbf e_t^\top\mid X_{t-1},X_t]=0.
\]
因此交叉项期望为 0。

噪声协方差假设
\[
  \E[\mathbf e_s\mathbf e_s^\top\mid X_s]=\sigma^2I
\]
表示：噪声平均为 0，每个方向的方差都是 \(\sigma^2\)，不同方向之间不相关。利用这个假设可得
\[
  \E[\norm{E_t^2e_{t-1}}^2\mid X_{t-1}]
  =
  \sigma^2\tr((E_t^2)^\top E_t^2),
\]
同理得到 \(E_t^3\) 那一项。两项相加就是公式 19。

\paragraph{第十步：为什么用 Markov 不等式，以及公式 20 怎么来。}
我们想给良性任务的检测分数找一个阈值，保证良性任务大概率不会被误报。\(d_t^2\) 是非负随机变量，所以最直接可用 Markov 不等式：
\[
  \Pr(X>a)\le\frac{\E[X]}{a}.
\]
令 \(X=d_t^2\)。希望单个时刻误报概率不超过 \(\epsilon/T\)，就取
\[
  a=\frac{T}{\epsilon}\E[d_t^2].
\]
也就是
\[
  \theta_t^2
  =
  \frac{T}{\epsilon}
  \sigma^2
  \left[
  \tr((E_t^2)^\top E_t^2)
  +
  \tr((E_t^3)^\top E_t^3)
  \right].
\]
开平方就得到公式 20。再对所有 \(t\) 用 union bound，就能保证整个训练过程中良性任务同时不过阈值的概率至少为 \(1-\epsilon\)。

\paragraph{第十一步：sub-Gaussian 和 Hanson--Wright 是什么。}
Markov 不等式很通用，但很松。若噪声是 sub-Gaussian，说明它的尾部像高斯分布一样下降很快，不容易出现极端大值。直观上：
\[
  \Pr(|X|>u)
\]
会随 \(u\) 增大快速下降。

Hanson--Wright 不等式专门控制 sub-Gaussian 随机向量的二次型：
\[
  \xi^\top A\xi.
\]
而这里
\[
  d_t^2
  =
  \xi_t^\top R_t^\top R_t\xi_t
\]
正是二次型。因此若噪声 sub-Gaussian，就可以用 Hanson--Wright 给出比 Markov 更紧的阈值，通常从 \(T/\epsilon\) 量级改善到 \(\log(T/\epsilon)\) 量级。

\paragraph{第十二步：正则化矩阵公式 21 从哪里来。}
T2T 最终收敛证明选用
\[
  H_t=\frac1{n_t}
  \left(
  \frac{\sigma^2}{W}I+\sum_{s=1}^{t-1}X_s^\top X_s
  \right).
\]
这个选择可以理解成“历史数据信息的累计”。\(\sum_{s=1}^{t-1}X_s^\top X_s\) 是过去任务给出的总曲率/信息量；\(\sigma^2/W\) 像一个初始先验强度，避免一开始矩阵不可逆；前面的 \(1/n_t\) 是为了和当前任务损失里的 \(X_t^\top X_t/n_t\) 尺度一致。

这个 \(H_t\) 的好处是递推会望远镜式简化，最终 \(w_T-w^\star\) 可以写成所有任务噪声和攻击的加权和，方便证明收敛。

\paragraph{第十三步：式 22 从哪里来。}
在线性回归中，带标签扰动的目标是
\[
  \frac1{2n_t}\norm{X_tw-(X_tw^\star+\mathbf e_t+\mathbf n_t)}^2
  +\frac12\norm{w-w_{t-1}}_{H_t}^2.
\]
对 \(w\) 求梯度并令 0：
\[
  \frac{X_t^\top}{n_t}
  \left[
  X_t(w-w^\star)-(\mathbf e_t+\mathbf n_t)
  \right]
  +H_t(w-w_{t-1})=0.
\]
把 \(w=w_t\)，并把含 \(w_t-w^\star\) 的项放左边：
\[
  \left(
  \frac{X_t^\top X_t}{n_t}+H_t
  \right)(w_t-w^\star)
  =
  \frac{X_t^\top(\mathbf e_t+\mathbf n_t)}{n_t}
  +H_t(w_{t-1}-w^\star).
\]
左乘逆矩阵，就得到式 22：
\[
  w_t-w^\star
  =
  \left(
  \frac{X_t^\top X_t}{n_t}+H_t
  \right)^{-1}
  \left[
  \frac{X_t^\top(\mathbf e_t+\mathbf n_t)}{n_t}
  +H_t(w_{t-1}-w^\star)
  \right].
\]
这说明最终递推仍然来自同一个原则：正则化训练目标求梯度等于 0。

\subsection{检测阈值：Lemma 4.3}

在线性回归假设下，如果任务 \(t-1\) 和 \(t\) 都是良性的，那么
\[
  d_t^2
  =
  \norm{E_t^2\mathbf e_{t-1}+E_t^3\mathbf e_t}_2^2 .
\]
噪声独立、零均值，且
\[
  \E[\mathbf e_s\mathbf e_s^\top\mid X_s]=\sigma^2I .
\]
因此交叉项期望为 0，
\begin{align}
  \E[d_t^2\mid X_{t-1},X_t]
  =
  \sigma^2\tr((E_t^2)^\top E_t^2)
  +
  \sigma^2\tr((E_t^3)^\top E_t^3).
  \label{eq:dt-expectation}
\end{align}
用 Markov 不等式：
\[
  \Pr(d_t^2>\theta_t^2)\le\epsilon/T,
\]
取
\begin{equation}
  \theta_t
  =
  \sqrt{
  \frac{\sigma^2T}{\epsilon}
  \left[
  \tr((E_t^2)^\top E_t^2)
  +
  \tr((E_t^3)^\top E_t^3)
  \right]
  }.
  \label{eq:theta}
\end{equation}
再对所有 \(t\) 用 union bound，得到同时成立概率至少 \(1-\epsilon\)。

如果噪声是 sub-Gaussian，可以用 Hanson-Wright 不等式替代 Markov 不等式，
阈值会从粗糙的 \(T/\epsilon\) 型改善成对数型 \(\log(T/\epsilon)\)。

\subsection{Theorem 4.4：最终误差界}

论文选用正则化矩阵
\begin{equation}
  H_t=\frac{1}{n_t}
  \left(
  \frac{\sigma^2}{W}I+\sum_{s=1}^{t-1}\mathbf X_s^\top\mathbf X_s
  \right).
  \label{eq:Ht-t2t}
\end{equation}
把攻击写成标签扰动 \(\mathbf n_t\)，其中未攻击时 \(\mathbf n_t=0\)。
线性回归递推为
\begin{align}
  \mathbf w_t-\mathbf w^\star
  =
  \left(
  \frac{\mathbf X_t^\top\mathbf X_t}{n_t}+H_t
  \right)^{-1}
  \left[
  \frac{\mathbf X_t^\top(\mathbf e_t+\mathbf n_t)}{n_t}
  +H_t(\mathbf w_{t-1}-\mathbf w^\star)
  \right].
  \label{eq:t2t-linear-recursion}
\end{align}
代入 \eqref{eq:Ht-t2t} 并展开递推，有
\begin{equation}
  \mathbf w_T-\mathbf w^\star
  =
  \tilde{\mathbf X}_{T+1}^{-1}
  \left[
  \sum_{s=1}^T\mathbf X_s^\top(\mathbf e_s+\mathbf n_s)
  +\frac{\sigma^2}{W}(\mathbf w_0-\mathbf w^\star)
  \right],
  \label{eq:t2t-unroll}
\end{equation}
其中
\[
  \tilde{\mathbf X}_{T+1}
  =
  \frac{\sigma^2}{W}I+\sum_{s=1}^T\mathbf X_s^\top\mathbf X_s .
\]

把任务分成良性任务和漏检攻击任务 \(\mathcal K=\{t_1,\dots,t_K\}\)。
检测到的攻击会被回滚丢弃，不进入最终递推；真正需要控制的是漏检攻击。
由 T2T 阈值可推出，任何漏检攻击必须满足
\begin{equation}
  \left\|
  \left(
  \frac{\mathbf X_t^\top\mathbf X_t}{n_t}
  \right)^+
  \frac{\mathbf X_t^\top}{n_t}
  (\mathbf n_t+\mathbf e_t)
  \right\|_2^2
  \le
  \frac{\sigma^2T}{\epsilon}
  \tr\!\left((\mathbf X_t^\top\mathbf X_t)^+\right).
  \label{eq:undetected-bound}
\end{equation}
直观解释：如果攻击超过这个尺度，在某个最坏未来任务对齐下，下一次 T2T 分数会超过阈值，攻击会被检测并回滚。

由 \eqref{eq:t2t-unroll}，
\begin{align}
  \E\norm{\mathbf w_T-\mathbf w^\star}^2
  \le&
  \sigma^2\tr(\tilde{\mathbf X}_{T+1}^{-1})
  \notag\\
  &+
  K\sum_{k=1}^K
  \left\|
  \tilde{\mathbf X}_{T+1}^{-1}
  (\mathbf X_{t_k}^\top\mathbf X_{t_k})
  \right\|_2^2
  \frac{\sigma^2T}{\epsilon}
  \tr\!\left((\mathbf X_{t_k}^\top\mathbf X_{t_k})^+\right).
  \label{eq:t2t-final-bound}
\end{align}
第一项是良性噪声误差；第二项是漏检攻击误差。

若所有任务 informative，即每个参数方向都能在 \(T\) 个任务中得到线性数量的信息，
则 \(\tilde{\mathbf X}_{T+1}\) 的特征值按 \(\Theta(T)\) 增长。
于是
\[
  \tr(\tilde{\mathbf X}_{T+1}^{-1})=O(1/T),
  \qquad
  \left\|
  \tilde{\mathbf X}_{T+1}^{-1}
  (\mathbf X_{t_k}^\top\mathbf X_{t_k})
  \right\|_2^2=O(1/T^2).
\]
代入 \eqref{eq:t2t-final-bound} 得
\[
  \calR(w_T)=O\!\left(\frac{K^2}{\epsilon T}\right).
\]
当 \(K=O(1)\) 时，最终误差收敛到 0。

\section{频繁有界零均值攻击：鲁棒特征防御}

\subsection{先不看公式：第二种防御的整体逻辑}

T2T 适合“攻击次数少但每次可能很强”的情况。它的策略像安检：发现可疑任务就丢掉或回滚。

但如果攻击很频繁，几乎每个任务都有一点污染，安检式防御就不适合了。因为你不能把大部分任务都丢掉，否则模型没有数据可学。所以第二种防御换了思路：不再主要判断“哪个任务坏了”，而是让模型在训练时自动降低对脆弱特征方向的敏感性。

可以把模型参数空间想象成很多方向：
\[
  u_1,u_2,\dots,u_p.
\]
有些方向数据很充分、噪声影响小；有些方向数据很弱、攻击稍微一推就会造成大误差。攻击者最喜欢打后者。

正则化矩阵 \(H_t\) 在这里就像一组方向刹车：
\begin{itemize}
  \item 某个方向的正则化强，模型就不容易沿这个方向被当前任务数据拉走。
  \item 某个方向的正则化弱，模型就更愿意吸收当前任务在这个方向的信息。
\end{itemize}
鲁棒特征防御要做的是：给每个方向选择合适的刹车强度，让攻击者即使知道模型，也很难找到一个特别脆弱的方向集中攻击。

这一节的符号可以这样理解：
\begin{itemize}
  \item \(u_j\)：第 \(j\) 个特征方向。
  \item \(\gamma_j^t\)：任务 \(t\) 在方向 \(u_j\) 上提供了多少数据信息。越大，说明这个方向数据越有力。
  \item \(\lambda_j^t\)：防御者在方向 \(u_j\) 上设置的正则化强度，也就是刹车强度。
  \item \(\calR_j^t\)：模型在方向 \(u_j\) 上还剩多少误差。
  \item \(M\)：攻击者总预算。
\end{itemize}

为什么要对角化？如果不对角化，各方向混在一起，一个方向的更新会影响另一个方向，分析非常复杂。对角化的目的不是炫技，而是把一个高维矩阵问题拆成许多一维问题：
\[
  \text{第 }1\text{ 个方向怎么设刹车，}
  \text{第 }2\text{ 个方向怎么设刹车，}\dots
\]
这样才能解释“攻击者会打哪个方向”和“防御者应该保护哪些方向”。

为什么会有 protected set？因为攻击者预算有限，不会平均打所有方向，而会集中攻击最敏感的一组方向。防御者识别出这组方向后，把它们作为 protected set，专门调整这些方向的正则化强度，避免最坏方向单独决定总误差。

\subsection{为什么不能继续靠检测}

如果几乎每个任务都被攻击，T2T 检测即使有效，也会不断丢任务，模型无法学习。
此时要换思路：不再拒绝任务，而是调整正则化矩阵 \(H_t\)，让模型对攻击最敏感的方向变得不敏感。

\subsection{线性模型中的 minimax 递推}

在向量化线性回归下，带 non-shifted 攻击的递推是
\begin{equation}
  w_t-w_\star
  =
  S_t^{-1}
  \left[
  \frac{X_t^\top(\varepsilon_t+\tilde\eta_t)}{n_t}
  +H_t(w_{t-1}-w_\star)
  \right],
  \qquad
  S_t=\frac{X_t^\top X_t}{n_t}+H_t .
  \label{eq:rfd-recursion}
\end{equation}
因为 \(\tilde\eta_t\) 是零均值并与历史误差条件不相关，
\begin{align}
  \E[w_t-w_\star]
  &=
  S_t^{-1}H_t\E[w_{t-1}-w_\star],
  \label{eq:rfd-mean}\\
  \Cov(w_t-w_\star)
  &=
  S_t^{-1}\frac{X_t^\top}{n_t}
  (\Cov(\tilde\eta_t)+\sigma^2I)
  \frac{X_t}{n_t}S_t^{-1}
  \notag\\
  &\quad+
  S_t^{-1}H_t\Cov(w_{t-1}-w_\star)H_t^\top S_t^{-1}.
  \label{eq:rfd-cov}
\end{align}
于是防御者在每轮求解
\begin{align}
  \min_{H_t}\max_{\Cov(\tilde\eta_t)\succeq0}\quad
  &
  \norm{S_t^{-1}H_t\E[w_{t-1}-w_\star]}_2^2
  \notag\\
  &+
  \tr\!\left(
  S_t^{-1}\frac{X_t^\top}{n_t}
  (\Cov(\tilde\eta_t)+\sigma^2I)
  \frac{X_t}{n_t}S_t^{-1}
  \right)
  \notag\\
  &+
  \tr\!\left(
  S_t^{-1}H_t\Cov(w_{t-1}-w_\star)H_t^\top S_t^{-1}
  \right)
  \notag\\
  \st&
  \tr(\Cov(\tilde\eta_t))\le M .
  \label{eq:rfd-minimax}
\end{align}

\subsection{对角化假设与每个方向的风险}

假设所有任务 Hessian 两两可交换：
\[
  X_s^\top X_sX_t^\top X_t=X_t^\top X_tX_s^\top X_s .
\]
则它们可以被同一个正交矩阵 \(U=[u_1,\dots,u_p]\) 同时对角化：
\[
  \frac{X_t^\top X_t}{n_t}
  =
  U\Gamma_tU^\top,\qquad
  \Gamma_t=\diag(\gamma_1^t,\dots,\gamma_p^t).
\]
只需考虑
\[
  H_t=U\Lambda_tU^\top,\qquad
  \Lambda_t=\diag(\lambda_1^t,\dots,\lambda_p^t).
\]
定义第 \(j\) 个方向的误差贡献
\[
  \calR_j^t=\E[(u_j^\top(w_t-w_\star))^2],
  \qquad
  \calR(w_t)=\sum_{j=1}^p\calR_j^t.
\]

投影到 \(u_j\) 方向，若 \(\gamma_j^t>0\)，有
\begin{equation}
  u_j^\top(w_t-w_\star)
  =
  \frac{1}{\lambda_j^t+\gamma_j^t}
  \left[
  \frac{\sqrt{\gamma_j^t}}{\sqrt{n_t}}(v_j^t)^\top(\varepsilon_t+\tilde\eta_t)
  +\lambda_j^tu_j^\top(w_{t-1}-w_\star)
  \right],
  \label{eq:per-direction}
\end{equation}
其中 \(v_j^t=X_tu_j/\sqrt{n_t\gamma_j^t}\) 且 \(\norm{v_j^t}=1\)。
由零均值和不相关性，
\begin{equation}
  \calR_j^t
  =
  \left(\frac{\lambda_j^t}{\lambda_j^t+\gamma_j^t}\right)^2
  \calR_j^{t-1}
  +
  \frac{
  \gamma_j^t\left(\sigma^2+(v_j^t)^\top\Cov(\tilde\eta_t)v_j^t\right)
  }{
  n_t(\lambda_j^t+\gamma_j^t)^2
  } .
  \label{eq:risk-per-direction}
\end{equation}

固定 \(\lambda_j^t\) 后，攻击者的目标对 \(\Cov(\tilde\eta_t)\) 是线性的。
所以最优攻击把全部预算 \(M\) 放到最敏感方向：
\[
  j^\star\in
  \argmax_j
  \frac{\gamma_j^t}{n_t(\lambda_j^t+\gamma_j^t)^2},
  \qquad
  \Cov(\tilde\eta_t)=Mv_{j^\star}^t(v_{j^\star}^t)^\top .
\]
于是防御问题等价于
\begin{align}
  \min_{\lambda_1^t,\dots,\lambda_p^t}\quad
  &
  \max_j
  \frac{M\gamma_j^t}{n_t(\lambda_j^t+\gamma_j^t)^2}
  +
  \sum_{j=1}^p
  \left[
  \left(\frac{\lambda_j^t}{\lambda_j^t+\gamma_j^t}\right)^2
  \calR_j^{t-1}
  +
  \frac{\gamma_j^t\sigma^2}{n_t(\lambda_j^t+\gamma_j^t)^2}
  \right].
  \label{eq:scalar-minimax}
\end{align}

\subsection{Proposition 5.2：最优正则化的水位结构}

令
\[
  \tilde\gamma_j^t=\gamma_j^tn_t,\qquad
  \calI_t^+=\{j:\tilde\gamma_j^t>0,\ \calR_j^{t-1}>0\}.
\]
对 \(j\in\calI_t^+\)，定义
\[
  b_j^t=
  \frac{\calR_j^{t-1}\tilde\gamma_j^t+\sigma^2}
  {\calR_j^{t-1}\sqrt{\tilde\gamma_j^t}} .
\]
按 \(b_j^t\) 从小到大排序。再定义
\[
  A_j^t=
  \frac{
  M+j\sigma^2+\sum_{k=1}^j\calR_k^{t-1}\tilde\gamma_k^t
  }{
  \sum_{k=1}^j\calR_k^{t-1}\sqrt{\tilde\gamma_k^t}
  } .
\]
令 \(j_t\) 是满足
\[
  A_j^t>b_j^t
\]
的最大 \(j\)。那么最优防御为
\begin{equation}
  \lambda_j^t=
  \begin{cases}
  \dfrac{\sigma^2}{n_t\calR_j^{t-1}},
  & j\in\calI_t^+,\ j>j_t,\\[0.9em]
  A_{j_t}^t\sqrt{\dfrac{\gamma_j^t}{n_t}}-\gamma_j^t,
  & j\in\calI_t^+,\ j\le j_t,\\[0.9em]
  +\infty,
  & \gamma_j^t>0,\ \calR_j^{t-1}=0,\\[0.4em]
  \text{任意正数},
  & \gamma_j^t=0 .
  \end{cases}
  \label{eq:rfd-lambda}
\end{equation}

\subsection{Proposition 5.2 的证明}

证明从变量替换开始：
\[
  c_j^t=\frac{\lambda_j^t}{\lambda_j^t+\gamma_j^t}.
\]
这表示第 \(j\) 个方向保留旧模型的比例；\(1-c_j^t\) 表示当前任务数据进入模型的比例。

忽略退化方向后，问题变成
\begin{align}
  \min_{\{c_j^t\},z}\quad
  &
  \sum_{j\in\calI_t^+}
  \left[
  (c_j^t)^2\calR_j^{t-1}
  +
  \frac{(1-c_j^t)^2\sigma^2}{n_t\gamma_j^t}
  \right]
  +zM
  \notag\\
  \st&
  z\ge
  \frac{(1-c_j^t)^2}{\gamma_j^tn_t},
  \qquad j\in\calI_t^+ .
  \label{eq:kkt-problem}
\end{align}
这里 \(z\) 代表最大攻击敏感度。

对约束引入乘子 \(\chi_j\ge0\)。KKT 条件给出
\[
  \sum_{j\in\calI_t^+}\chi_j=M
\]
以及
\begin{equation}
  c_j^t=
  \frac{(\sigma^2+\chi_j)/(n_t\gamma_j^t)}
  {\calR_j^{t-1}+(\sigma^2+\chi_j)/(n_t\gamma_j^t)} .
  \label{eq:c-kkt}
\end{equation}
若 \(\chi_j>0\)，说明方向 \(j\) 的最大敏感度约束是活跃的：
\[
  z=\frac{(1-c_j^t)^2}{\gamma_j^tn_t}.
\]
把 \eqref{eq:c-kkt} 代入可得
\begin{equation}
  \chi_j
  =
  \frac{\calR_j^{t-1}\sqrt{\gamma_j^tn_t}}{\sqrt z}
  -\sigma^2
  -\gamma_j^tn_t\calR_j^{t-1}.
  \label{eq:chi}
\end{equation}

设活跃集合为 \(J_t=\{j:\chi_j>0\}\)。由 \(\sum\chi_j=M\) 得
\[
  \frac{1}{\sqrt z}
  =
  \frac{
  M+|J_t|\sigma^2+\sum_{j\in J_t}\gamma_j^tn_t\calR_j^{t-1}
  }{
  \sum_{j\in J_t}\calR_j^{t-1}\sqrt{\gamma_j^tn_t}
  } .
\]
这就是 \(A_j^t\) 的来源。

接下来要确定 \(J_t\)。令
\[
  a_j=\calR_j^{t-1}\sqrt{\gamma_j^tn_t},
  \qquad
  b_j=
  \frac{\gamma_j^tn_t\calR_j^{t-1}+\sigma^2}
  {\calR_j^{t-1}\sqrt{\gamma_j^tn_t}} .
\]
对任意候选活跃集合 \(J\)，KKT 条件给出
\[
  \chi_k=a_k(A_J-b_k).
\]
因此：
\[
  k\in J \Rightarrow b_k<A_J,
  \qquad
  k\notin J \Rightarrow b_k\ge A_J .
\]
所以活跃集合必须是按 \(b_j\) 排序后的一个前缀。
最大满足 \(A_j^t>b_j^t\) 的 \(j\) 就是前缀长度 \(j_t\)。

最后由
\[
  \lambda_j^t=\frac{c_j^t\gamma_j^t}{1-c_j^t}
\]
还原回 \(\lambda_j^t\)，得到 \eqref{eq:rfd-lambda}。

\begin{remark}
\(b_j^t\) 越小，方向 \(j\) 越容易进入 protected set。直观上，这些方向要么当前特征信息弱，要么历史误差大，要么噪声相对影响强，因此更容易被攻击者利用。
\end{remark}

\subsection{Theorem 5.3：为什么鲁棒特征防御更快}

定理考虑一个特殊平稳情形：protected set 长期固定为
\[
  \{1,\dots,j_T\}.
\]
对 protected directions 定义总误差
\[
  \Phi_t=\sum_{j=1}^{j_T}\calR_j^t,
  \qquad
  f_j^t=\frac{\calR_j^t}{\Phi_t}.
\]
假设 \(f_j^t\to f_j\)，且 \(\tilde\gamma_j^t\to\tilde\gamma_j^T\)。

在 protected directions 上，KKT 条件让攻击敏感度相等：
\[
  \frac{(1-c_j^t)^2}{\tilde\gamma_j^t}=z_t,
  \qquad j=1,\dots,j_T.
\]
因此
\[
  c_j^t=1-\sqrt{z_t\tilde\gamma_j^t}.
\]
由 KKT 水位公式：
\[
  \sqrt{z_t}
  =
  \frac{
  \sum_{i=1}^{j_T}\sqrt{\tilde\gamma_i^t}\calR_i^{t-1}
  }{
  M+j_T\sigma^2+\sum_{i=1}^{j_T}\tilde\gamma_i^t\calR_i^{t-1}
  } .
\]

把 protected directions 的风险加总，得到核心递推：
\begin{equation}
  \Phi_t
  =
  \Phi_{t-1}
  -
  \frac{
  \left(
  \sum_{j=1}^{j_T}\sqrt{\tilde\gamma_j^t}\calR_j^{t-1}
  \right)^2
  }{
  M+j_T\sigma^2+
  \sum_{j=1}^{j_T}\tilde\gamma_j^t\calR_j^{t-1}
  } .
  \label{eq:phi-recursion}
\end{equation}
把 \(\calR_j^{t-1}=f_j^{t-1}\Phi_{t-1}\) 代入，并取倒数，可写成
\begin{align}
  \frac{1}{\Phi_t}-\frac{1}{\Phi_{t-1}}
  =
  \frac{
  \left(
  \sum_{j=1}^{j_T}f_j^{t-1}\sqrt{\tilde\gamma_j^t}
  \right)^2
  }{
  M+j_T\sigma^2+
  \Phi_{t-1}
  \left[
  \sum_{j=1}^{j_T}f_j^{t-1}\tilde\gamma_j^t
  -
  \left(
  \sum_{j=1}^{j_T}f_j^{t-1}\sqrt{\tilde\gamma_j^t}
  \right)^2
  \right]
  }.
  \label{eq:phi-reciprocal}
\end{align}
Cauchy 不等式保证中括号非负，因此 \(\Phi_t\) 单调下降。
在平稳假设下，分子收敛到
\[
  C_\star=
  \left(
  \sum_{j=1}^{j_T}f_j\sqrt{\tilde\gamma_j^T}
  \right)^2>0,
\]
而 \(\Phi_t\to0\)，所以分母最终趋于 \(M+j_T\sigma^2\)。
于是
\[
  \frac{1}{\Phi_t}-\frac{1}{\Phi_{t-1}}
  \to
  \frac{C_\star}{M+j_T\sigma^2}.
\]
累加并用 Cesaro 平均得到
\[
  \Phi_T
  \sim
  \frac{M+j_T\sigma^2}
  {T\left(\sum_{j=1}^{j_T}f_j\sqrt{\tilde\gamma_j^T}\right)^2}.
\]

对非 protected directions，攻击者不投预算，递推退化为
\[
  \calR_j^t
  =
  \frac{\sigma^2\calR_j^{t-1}}
  {\sigma^2+\tilde\gamma_j^t\calR_j^{t-1}}.
\]
取倒数：
\[
  \frac{1}{\calR_j^t}
  =
  \frac{1}{\calR_j^{t-1}}
  +\frac{\tilde\gamma_j^t}{\sigma^2}.
\]
展开得
\[
  \calR_j^T
  \le
  \frac{\sigma^2}
  {\sigma^2/W+\sum_{t=1}^T\tilde\gamma_j^t}.
\]
因此鲁棒防御的总误差为
\begin{equation}
  \calR(w_T)
  \lesssim
  \frac{M+j_T\sigma^2}
  {T\left(\sum_{j=1}^{j_T}f_j\sqrt{\tilde\gamma_j^T}\right)^2}
  +
  \sum_{j=j_T+1}^p
  \frac{\sigma^2}
  {\sigma^2/W+\sum_{\tau=1}^T\tilde\gamma_j^\tau}.
  \label{eq:rfd-final}
\end{equation}

\subsection{和普通正则化 / EWC 型基线的比较}

基线使用
\[
  H_t=\frac{1}{n_t}
  \left(
  \frac{\sigma^2}{W}I+\sum_{\tau=1}^{t-1}X_\tau^\top X_\tau
  \right).
\]
在方向 \(j\) 上，最终噪声权重约为
\[
  \frac{\sqrt{\tilde\gamma_j^t}}
  {\sigma^2/W+\sum_{\tau=1}^T\tilde\gamma_j^\tau}.
\]
攻击者每轮把预算投到最坏方向，攻击项变成
\[
  M\max_j
  \frac{\sum_{t=1}^T\tilde\gamma_j^t}
  {\left(\sigma^2/W+\sum_{t=1}^T\tilde\gamma_j^t\right)^2}.
\]
平稳时
\[
  \sum_{t=1}^T\tilde\gamma_j^t\sim T\tilde\gamma_j^T,
\]
于是基线攻击项近似
\[
  \max_j\frac{M}{T\tilde\gamma_j^T}.
\]
这意味着一个小 \(\tilde\gamma_j^T\) 的脆弱方向就可能主导误差。
鲁棒特征防御则用 protected set 把敏感度平均化，攻击项变为
\[
  \frac{M+j_T\sigma^2}
  {T\left(\sum_{j=1}^{j_T}f_j\sqrt{\tilde\gamma_j^T}\right)^2},
\]
不再完全由单个最坏方向决定。

\section{实验部分怎么读}

实验不是证明的核心，但它验证了两个理论直觉。

\begin{itemize}
  \item 对 infrequent shifted attacks，T2T 检测分数 \(d_t\) 在攻击任务附近出现峰值。即使测试准确率看起来没明显变化，微观更新分数仍能暴露攻击。
  \item 对 frequent bounded non-shifted attacks，iCaRL 这类 replay/memorization 方法在频繁攻击下可能更脆弱，因为它没有显式正则化矩阵来降低敏感方向。鲁棒特征防御比普通 EWC 型正则化收敛更快。
\end{itemize}

\section{读公式的路线图}

如果从头读论文很吃力，可以按以下顺序理解。

\begin{enumerate}
  \item 先看 \eqref{eq:regularized-cl}。所有防御都只是在设计 \(H_t\) 或通过检测让 \(H_t\to\infty\)。
  \item 再看 \eqref{eq:attack-surrogate}。它告诉你攻击者为什么打最大奇异方向。
  \item 用一维递推 \eqref{eq:scalar-recursion} 理解 Theorem 3.1。这个定理不需要复杂矩阵。
  \item 用等式 \(D_t^1A_t=D_t^2B_t\) 理解 T2T。它解释了为什么 \(d_t\) 能消去历史。
  \item 用 \eqref{eq:t2t-unroll} 理解 T2T 的收敛界。最终模型就是所有噪声与攻击的加权和。
  \item 用 \eqref{eq:risk-per-direction} 理解鲁棒特征防御。每个特征方向有一个独立风险递推。
  \item 用 KKT 条件理解 \eqref{eq:rfd-lambda}。protected set 是一个排序后的前缀集合。
  \item 最后看 \eqref{eq:phi-reciprocal}。Theorem 5.3 的 \(O(1/T)\) 来自倒数递推。
\end{enumerate}

\section{这篇论文的假设与局限}

\begin{itemize}
  \item 最完整的证明建立在线性回归和平方损失上。论文用 NTK 和预训练线性头解释为什么这对非线性模型有启发，但严格结论仍然依赖线性化结构。
  \item T2T 防御需要维护最近几个模型、正则化矩阵和 Hessian 信息。高维神经网络中必须用 Hessian 近似。
  \item Theorem 4.4 中 informative tasks 假设很关键。若某些参数方向长期没有数据信息，即使没有攻击也不能保证收敛。
  \item Theorem 5.3 的闭式收敛率依赖 protected set 平稳、Hessian 可同时对角化等假设。这些假设主要用于给出可解释的理论形式。
  \item 攻击预算 \(M\) 在鲁棒特征防御中需要被知道或至少有上界。若上界过松，正则化可能过保守。
\end{itemize}

\section{一句话总结}

这篇论文的核心观点是：连续学习中的投毒攻击不是普通噪声，而是沿时间累积的策略性偏置。频繁均值偏移或巨大方差攻击会打破任何正则化 CL 的收敛保证；少量强攻击可以通过 task-to-task 投影检测并回滚；频繁有界零均值攻击则需要在 Hessian 特征方向上做 minimax 正则化，把攻击者最想利用的敏感方向压平。

\begin{thebibliography}{9}
\bibitem{hu2026}
Yiting Hu and Lingjie Duan.
\newblock \emph{Theory of Continual Learning Against Data Poisoning Attacks}.
\newblock arXiv:2606.29841, 2026.
\end{thebibliography}

\end{document}
